% ============================================================
%  ch02_framework.tex — Mathematical Framework
%  Requires: macros.tex, amsmath, amssymb, amsthm
%  Estimated length: 25–35 pages
%  Rewritten: VD-08, 2026-03-07
% ============================================================

\chapter{Mathematical Framework}\label{ch:framework}

Our goal in this chapter is to build, from first principles, the
mathematical machinery required for the departure-variable approach
to Bianchi anisotropy bounds.  We begin with the 1+3 covariant
decomposition of spacetime in its most general form---retaining
vorticity and four-acceleration throughout---derive the
propagation and constraint equations from the Ricci identity,
introduce standardised kinematic variables and their
normalisation convention, construct FLRW comparator geometries,
prove the master defect identity, set up the three-frame problem
and the King--Ellis frame-attribution corrections, and close with
the tilted Raychaudhuri equation expressed in departure variables.
Every derivation that later chapters rely on is given in full;
results that serve only as context are stated with references.

For the revision programme, this chapter also fixes the ownership
boundary of the claim-tiered framework.  The scalar coordinates
$x$, $Q$, $\Pi$, $F$, and $G_F$ are diagnostic or likelihood inputs
only after their owner, transfer source, null status, covariance status,
and caveats are declared.  MIO may report diagnostic certificates and
direction/depth coherence; HTT owns model-dependent likelihoods,
posterior pushforwards, and evidence surfaces.  Neither layer turns
scalar compatibility into geometry detection or Bianchi family
identification.

This is the formal no-go anchor for the rearchitected
manuscript.  A scalar-amplitude compression can support a
claim-tiered observational/statistical framework, a transfer
conditional diagnostic, or a pre-inference rank/null gate, but it
cannot by itself identify a Bianchi geometry, select a Bianchi
family, or certify a global-tilt source.  Family-level promotion
requires the future native low-ell morphology atlas, matched sky
support, covariance/null calibration, and equivalence-class gates.


% ============================================================
\section{The 1+3 covariant decomposition}\label{sec:13cov}
% ============================================================

\subsection{Choice of timelike congruence}\label{sec:congruence}

Before diving into the algebraic machinery, we clarify the
starting point.  General relativity provides no preferred time
direction: the spacetime manifold $(\mathcal{M}, g_{ab})$ is a
four-dimensional Lorentzian manifold, and any timelike vector
field may serve as a local ``clock.''  In cosmology, a natural
choice arises from the matter content.  A perfect fluid with
four-velocity $\hat{u}^a$ picks out a preferred congruence of
worldlines; a spatially homogeneous geometry picks out a
group-invariant time direction $u^a$.  These two choices coincide
(the fluid is \emph{comoving}) if and only if the tilt rapidity
$\beta = 0$.  When tilt is non-zero, the distinction between the
\emph{geometry frame} $u^a$ and the \emph{matter frame}
$\hat{u}^a$ is central to the analysis
(Section~\ref{sec:frame-problem}).

For the 1+3 decomposition, we work with an arbitrary unit
timelike vector field $u^a$ satisfying
\begin{equation}\label{eq:u-norm}
u_a\, u^a = -1\,.
\end{equation}
The minus sign follows from the $(-,+,+,+)$ metric signature.
The four-velocity $u^a$ defines, at each point, a rest space
orthogonal to itself, and we now construct the projection
operator onto that space.


\subsection{The spatial projector}\label{sec:projector}

\begin{definition}[Spatial projector]\label{def:hab}
The tensor
\begin{equation}\label{eq:hab-def}
h_{ab} := g_{ab} + u_a\, u_b
\end{equation}
projects any spacetime tensor onto the local rest space of $u^a$.
\end{definition}

We verify the essential algebraic properties.

\begin{proposition}[Properties of $h_{ab}$]
\label{prop:hab-properties}
The projector $h_{ab}$ satisfies:
\begin{enumerate}
\item[\emph{(i)}] \textbf{Orthogonality:} $h_{ab}\,u^b = 0$.
\item[\emph{(ii)}] \textbf{Idempotency:}
  $h^a{}_c\,h^c{}_b = h^a{}_b$.
\item[\emph{(iii)}] \textbf{Trace:} $h^a{}_a = 3$.
\item[\emph{(iv)}] \textbf{Symmetry:} $h_{ab} = h_{ba}$.
\end{enumerate}
\end{proposition}

\begin{proof}
(i)~Direct computation:
$h_{ab}\,u^b = g_{ab}\,u^b + u_a\,u_b\,u^b
= u_a + u_a(-1) = 0$.

(ii)~Contract:
$h^a{}_c\,h^c{}_b
= (\delta^a_c + u^a u_c)(\delta^c_b + u^c u_b)
= \delta^a_b + u^a u_b + u^a u_b + u^a(u_c u^c)u_b
= \delta^a_b + 2u^a u_b + u^a(-1)u_b
= \delta^a_b + u^a u_b = h^a{}_b$.

(iii)~$h^a{}_a = \delta^a_a + u^a u_a = 4 + (-1) = 3$.

(iv)~Both $g_{ab}$ and $u_a u_b$ are symmetric.
\end{proof}

We call a tensor \emph{spatial} (with respect to $u^a$) if
every free index, upon contraction with $u^a$, gives zero.  The
projector $h_{ab}$ extracts the spatial part of any tensor.  For
instance, for a vector $V^a$, the spatial part is
$h^a{}_b V^b = V^a + u^a(u_b V^b)$.  A spatial vector has
three independent components---the rest space is
three-dimensional, consistent with $h^a{}_a = 3$.


\subsection{Irreducible decomposition of $\nabla_b u_a$}
\label{sec:irred-decomp}

The covariant derivative $\nabla_b u_a$ is a
$(0,2)$~spacetime tensor with $4 \times 4 = 16$ independent
components.  We decompose it into parts that are spatial (in
both indices) and a part that involves the time direction $u^b$.
The decomposition proceeds in three steps.

\paragraph{Step 1: Time--space split.}
Project both indices of $\nabla_b u_a$ using $h$ and $u$:
\begin{equation}\label{eq:ts-split}
\nabla_b u_a
= h^c{}_a\, h^d{}_b\, \nabla_d u_c
  + h^c{}_a\, u^d u_b\, \nabla_d u_c
  + u^c u_a\, h^d{}_b\, \nabla_d u_c
  + u^c u_a\, u^d u_b\, \nabla_d u_c\,.
\end{equation}
Now we evaluate each term.  The quantity
$u^d \nabla_d u_c =: \dot{u}_c$ is the four-acceleration.
Since $u_c u^c = -1$, differentiating gives
$u^c \dot{u}_c = 0$---the acceleration is spatial.  The
last two terms in Eq.~\eqref{eq:ts-split} involve
$u^c u_a \,(\cdots)_c$; since $\nabla_d u_c$ contracted with
$u^c$ gives $\frac{1}{2}\nabla_d(u_c u^c) = 0$, the third and
fourth terms both vanish.  The second term gives
$- \dot{u}_a\, u_b$ (using $h^c{}_a \dot{u}_c = \dot{u}_a$
since $\dot{u}_a$ is spatial).  Hence
\begin{equation}\label{eq:ts-result}
\nabla_b u_a
= h^c{}_a\, h^d{}_b\, \nabla_d u_c
  - \dot{u}_a\, u_b\,.
\end{equation}
The first term is the fully spatial part of $\nabla_b u_a$; we
denote it $\tilde{\nabla}_b u_a := h^c{}_a h^d{}_b \nabla_d u_c$.
This spatial tensor has $3 \times 3 = 9$ independent components.

\paragraph{Step 2: Symmetric--antisymmetric split.}
Decompose $\tilde{\nabla}_b u_a$ into its symmetric and
antisymmetric parts:
\begin{equation}
\tilde{\nabla}_b u_a
= \tilde{\nabla}_{(b} u_{a)}
  + \tilde{\nabla}_{[b} u_{a]}\,.
\end{equation}
The symmetric part has $6$ components; the antisymmetric part
has $3$.

\paragraph{Step 3: Trace extraction.}
The symmetric part decomposes further into a trace and a
trace-free piece.  The trace is
\begin{equation}
h^{ab}\tilde{\nabla}_b u_a
= h^{ab}\nabla_b u_a
= \nabla_a u^a + u^a u^b \nabla_b u_a
= \nabla_a u^a + u^a \dot{u}_a
= \nabla_a u^a =: \Theta\,,
\end{equation}
where we used $u^a \dot{u}_a = 0$.  The scalar $\Theta$ is the
\emph{volume expansion}.  Removing the trace from the symmetric
part gives the \emph{shear tensor}
\begin{equation}\label{eq:sigma-extract}
\sigma_{ab}
:= \tilde{\nabla}_{(b} u_{a)}
   - \tfrac{1}{3}\,\Theta\, h_{ab}\,,
\end{equation}
which is symmetric ($\sigma_{ab} = \sigma_{(ab)}$), spatial
($\sigma_{ab} u^b = 0$), and trace-free
($\sigma^a{}_a = 0$) by construction.  The antisymmetric part
defines the \emph{vorticity tensor}
\begin{equation}\label{eq:omega-extract}
\omega_{ab}
:= \tilde{\nabla}_{[b} u_{a]}\,,
\end{equation}
which is antisymmetric ($\omega_{ab} = -\omega_{ba}$) and spatial
($\omega_{ab} u^b = 0$).

Collecting everything, we arrive at the \emph{kinematic
decomposition}:
\begin{equation}\label{eq:u-decomp}
\boxed{
\nabla_b u_a
= \tfrac{1}{3}\,\Theta\,h_{ab}
  + \sigma_{ab} + \omega_{ab}
  - \dot{u}_a\, u_b
}\,.
\end{equation}
This identity is exact and holds for any unit timelike
congruence in any spacetime.  The $16$ components of
$\nabla_b u_a$ split as:
$1\;(\Theta) + 5\;(\sigma_{ab}) + 3\;(\omega_{ab})
+ 3\;(\dot{u}_a) + 4\;\text{(identically zero by constraints)}
= 16$.

\begin{definition}[PSTF projection]\label{def:pstf}
For any spatial rank-2 tensor $A_{ab}$, the \emph{projected
symmetric trace-free} (PSTF) part is
\begin{equation}\label{eq:pstf-def}
A_{\langle ab\rangle}
:= h^c{}_{(a}\,h^d{}_{b)}\,A_{cd}
   - \tfrac{1}{3}\,h_{ab}\,h^{cd}A_{cd}\,.
\end{equation}
Angle brackets on indices denote this projection. In this
notation, $\sigma_{ab} = \nabla_{\langle b} u_{a\rangle}$.
\end{definition}


\subsection{Physical interpretation of the kinematic
  quantities}\label{sec:phys-interp}

The \emph{expansion} $\Theta = \nabla_a u^a$ measures the
rate of change of the volume element carried by the congruence.
For a small fluid element with three-volume $V$, we have
$\dot{V}/V = \Theta$, so positive $\Theta$ means expansion and
negative $\Theta$ means contraction.  The Hubble rate is
$\Htheta := \Theta/3$.

The \emph{shear} $\sigma_{ab}$ measures the rate at which the
shape of a fluid element distorts at constant volume.  Because
$\sigma_{ab}$ is trace-free, it encodes distortion without
volume change; because it is symmetric, it has five independent
components (three eigenvalues summing to zero and two Euler
angles specifying the principal axes).  In an FLRW spacetime,
$\sigma_{ab} = 0$ identically.

The \emph{vorticity} $\omega_{ab}$ measures the local rotation
of the congruence.  The vorticity vector
$\omega^a := \frac{1}{2}\eta^{abcd}u_b\omega_{cd}$ (where
$\eta^{abcd}$ is the spacetime volume form) points along the
instantaneous rotation axis.  Non-zero vorticity prevents the
congruence from being hypersurface-orthogonal
(Section~\ref{sec:frobenius}).

The \emph{four-acceleration} $\dot{u}_a = u^b\nabla_b u_a$
measures the deviation of the worldlines from geodesic motion.
For a perfect barotropic fluid in a spatially homogeneous,
non-tilted model, $\dot{u}_a = 0$---the fluid elements follow
geodesics.  When tilt is present ($\beta \neq 0$) and
$p \neq 0$, King \& Ellis~\cite{KingEllis1973} proved that
$\dot{u}_a \neq 0$ necessarily: pressure gradients in the
tilted frame force the fluid off geodesics.

The scalar invariants of these tensors are
\begin{equation}\label{eq:scalar-invariants}
|\sigma|^2 := \tfrac{1}{2}\sigma_{ab}\sigma^{ab}\,,\qquad
|\omega|^2 := \tfrac{1}{2}\omega_{ab}\omega^{ab}\,,\qquad
|\dot{u}|^2 := \dot{u}_a\dot{u}^a\,.
\end{equation}
Each is non-negative by construction and vanishes if and only if
the corresponding tensor vanishes.


\subsection{Energy--momentum decomposition}\label{sec:em-decomp}

The energy--momentum tensor of a general fluid, decomposed
relative to $u^a$, takes the form
\begin{equation}\label{eq:Tab-general}
T_{ab} = \rho\, u_a u_b + p\, h_{ab}
  + 2\,q_{(a}\,u_{b)} + \pi_{ab}\,,
\end{equation}
where
$\rho = T_{ab}\,u^a u^b$ is the energy density,
$p = \frac{1}{3} h^{ab} T_{ab}$ is the isotropic pressure,
$q_a = -h^c{}_a\,T_{cd}\,u^d$ is the energy flux (heat flow),
and $\pi_{ab}$ is the PSTF anisotropic stress.
For a \emph{perfect fluid} observed in its rest frame
$\hat{u}^a$, we have $q_a = 0$ and $\pi_{ab} = 0$:
\begin{equation}\label{eq:Tab-pf}
T_{ab} = (\hat{\rho} + \hat{p})\,\hat{u}_a\hat{u}_b
       + \hat{p}\,g_{ab}\,.
\end{equation}
When the matter four-velocity $\hat{u}^a$ is tilted relative
to the geometry-frame congruence $u^a$ by a rapidity $\beta$,
the Lorentz boost gives
\begin{equation}\label{eq:boost-u}
\hat{u}^a = \cosh\beta\;(u^a + v^a)\,,
\end{equation}
where $v^a$ is spatial ($v_a u^a = 0$) with
$|v| = \tanh\beta$.  Projecting Eq.~\eqref{eq:Tab-pf} onto
the geometry frame, the effective energy density becomes
\begin{equation}\label{eq:rho-geom}
\rho_{\mathrm{tot}} := T_{ab}\,u^a u^b
= \hat{\rho} + (\hat{\rho} + \hat{p})\sinh^2\!\beta\,.
\end{equation}
The tilt contributes a positive energy density
$(\hat{\rho} + \hat{p})\sinh^2\!\beta$ to the geometry-frame
energy budget, sourced by the relativistic kinetic energy of
the fluid's bulk motion.


\subsection{The Raychaudhuri equation}\label{sec:raychaudhuri}

We now derive the evolution equation for the expansion scalar
from first principles.  The starting point is the Ricci identity
for the vector field $u^a$:
\begin{equation}\label{eq:ricci-identity}
(\nabla_a \nabla_b - \nabla_b \nabla_a)\,u^c
= R^c{}_{dab}\,u^d\,,
\end{equation}
where $R^c{}_{dab}$ is the Riemann tensor.  We contract on $c$
and $a$:
\begin{equation}\label{eq:ricci-contract1}
\nabla_a \nabla_b u^a - \nabla_b \nabla_a u^a
= R_{db}\,u^d\,,
\end{equation}
and then contract with $u^b$:
\begin{equation}\label{eq:raych-step1}
u^b \nabla_a \nabla_b u^a
- u^b \nabla_b (\nabla_a u^a)
= R_{db}\,u^d u^b\,.
\end{equation}
The second term on the left is $-\dot{\Theta}$ (the covariant
time derivative of the expansion).  The first term requires more
work.  Using the product rule,
\begin{equation}
u^b \nabla_a \nabla_b u^a
= \nabla_a(u^b \nabla_b u^a) - (\nabla_a u^b)(\nabla_b u^a)\,.
\end{equation}
The first piece is $\nabla_a \dot{u}^a$.  For the second, we
substitute the kinematic decomposition~\eqref{eq:u-decomp}:
\begin{equation}
(\nabla_a u^b)(\nabla_b u^a)
= \bigl(\tfrac{1}{3}\Theta\,h^b{}_a + \sigma^b{}_a
  + \omega^b{}_a - \dot{u}^b u_a\bigr)
  \bigl(\tfrac{1}{3}\Theta\,h_{b}{}^{a} + \sigma_{b}{}^{a}
  + \omega_{b}{}^{a} - \dot{u}_b u^a\bigr)\,.
\end{equation}
We expand this product and use the orthogonality and trace
properties: $h^b{}_a h_b{}^a = 3$,
$\sigma^b{}_a \sigma_b{}^a = 2|\sigma|^2$,
$\omega^b{}_a \omega_b{}^a = -2|\omega|^2$ (the minus sign arises
because $\omega_{ab}$ is antisymmetric), the cross terms between
$h_{ab}$ and the trace-free tensors vanish, and terms involving
$u_a \dot{u}^a = 0$ drop out.  The result is
\begin{equation}
(\nabla_a u^b)(\nabla_b u^a)
= \tfrac{1}{3}\Theta^2 + 2|\sigma|^2 - 2|\omega|^2\,.
\end{equation}
(Note that the ``time $\times$ time'' cross term
$(-\dot{u}^b u_a)(-\dot{u}^a u_b)
= (\dot{u}^a u_a)(\dot{u}^b u_b) = 0$
vanishes because $\dot{u}^a u_a = 0$.)
Assembling everything and introducing the Einstein field
equation $R_{ab} - \frac{1}{2}R\,g_{ab} + \Lambda\,g_{ab}
= \kappa\,T_{ab}$, which gives
$R_{ab}\,u^a u^b = \frac{1}{2}\kappa(\rho + 3p) - \Lambda$
(for a perfect fluid), we obtain
\begin{equation}
\dot{\Theta}
= -\tfrac{1}{3}\Theta^2
  - 2(|\sigma|^2 - |\omega|^2)
  - \tfrac{1}{2}\kappa(\rho + 3p)
  + \Lambda
  + \nabla_a\dot{u}^a\,.
\end{equation}
For a spatial vector $V^a$ (with $V_a u^a = 0$), the full and
projected divergences are related by
$\nabla_a V^a = \tilde{\nabla}_a V^a + \dot{u}_a V^a$.
Applying this identity to $V^a = \dot{u}^a$ gives
$\nabla_a\dot{u}^a = \tilde{\nabla}_a\dot{u}^a + |\dot{u}|^2$.
The \textbf{Raychaudhuri equation} thus takes its standard form:
\begin{equation}\label{eq:raychaudhuri}
\boxed{
\dot{\Theta}
= -\tfrac{1}{3}\Theta^2
  - 2(|\sigma|^2 - |\omega|^2)
  - \tfrac{1}{2}\kappa(\rho + 3p)
  + \Lambda
  + \tilde{\nabla}_a\dot{u}^a
  + \dot{u}_a\dot{u}^a
}\,.
\end{equation}
Note the sign structure: shear ($|\sigma|^2$) and matter
($\rho + 3p > 0$) drive $\Theta$ towards negative values
(focusing); vorticity ($|\omega|^2$), the cosmological constant
($\Lambda > 0$), and acceleration
($\tilde{\nabla}_a \dot{u}^a + |\dot{u}|^2$) resist focusing.  For
a geodesic irrotational congruence ($\dot{u}_a = 0$,
$\omega_{ab} = 0$) with $\Lambda \leq 0$ and
$\rho + 3p > 0$, the Raychaudhuri equation guarantees
$\dot{\Theta} < 0$: a singularity is unavoidable.


\subsection{Shear propagation equation}\label{sec:shear-prop}

We extract the PSTF part of the Ricci identity.  Starting
again from Eq.~\eqref{eq:ricci-identity}, we project both
free indices spatially and take the symmetric trace-free part.
A tedious and mundane computation (carrying through the kinematic
decomposition and collecting all PSTF-projected terms) yields
\begin{equation}\label{eq:shear-prop}
\dot{\sigma}_{\langle ab\rangle}
= -\tfrac{2}{3}\Theta\,\sigma_{ab}
  - \sigma_{c\langle a}\sigma^c{}_{b\rangle}
  + \omega_{\langle a}\omega_{b\rangle}
  - E_{ab}
  + \tfrac{1}{2}\kappa\pi_{ab}
  + \dot{u}_{\langle a}\dot{u}_{b\rangle}\,,
\end{equation}
where $E_{ab} = C_{acbd}\,u^c u^d$ is the electric part of the
Weyl tensor $C_{abcd}$, and $\pi_{ab}$ is the anisotropic
stress~\cite{EllisVanElst1999,EllisMaartens2012}.  The key
features of this equation are: (a)~the $-\frac{2}{3}\Theta\sigma_{ab}$
term provides Hubble damping of the shear; (b)~the Weyl tensor
$E_{ab}$ acts as a source, encoding tidal gravitational fields
from the anisotropic geometry; (c)~the acceleration term
$\dot{u}_{\langle a}\dot{u}_{b\rangle}$ provides an anisotropic
source from non-geodesic motion.


\subsection{Vorticity propagation equation}
\label{sec:vort-prop}

The antisymmetric part of the spatially projected Ricci
identity gives the vorticity propagation:
\begin{equation}\label{eq:vort-prop}
\dot{\omega}_{\langle ab\rangle}
= -\tfrac{2}{3}\Theta\,\omega_{ab}
  + \sigma_{c[a}\omega^c{}_{b]}
  + \tfrac{1}{2}\bigl(\nabla_{[a}\dot{u}_{b]}
    - u_{[a}\dot{\dot{u}}_{b]}\bigr)\,,
\end{equation}
where the last term involves the curl of the
acceleration~\cite{EllisVanElst1999}.  For a barotropic perfect
fluid, $\dot{u}_a = -\tilde{\nabla}_a p/(\rho+p)$, and the
curl of $\dot{u}_a$ vanishes identically---the relativistic
Kelvin theorem guarantees that vorticity-free initial data remain
vorticity-free.  For non-barotropic or multi-component fluids,
however, the curl of $\dot{u}_a$ can be non-zero, providing a
source for vorticity generation.  The shear--vorticity coupling
$\sigma_{c[a}\omega^c{}_{b]}$ transfers energy between shear
and vorticity modes without external sources.


% ============================================================
\section{Spatial curvature}\label{sec:spatial-curv}
% ============================================================

\subsection{The irrotational case and the Frobenius theorem}
\label{sec:frobenius}

When $\omega_{ab} = 0$, the congruence $u^a$ is
\emph{hypersurface-orthogonal}: there exist spacelike
hypersurfaces $\Sigma_t$ everywhere orthogonal to $u^a$.  The
Frobenius theorem provides the precise condition.

\begin{theorem}[Frobenius]\label{thm:frobenius}
A vector field $u^a$ is hypersurface-orthogonal if and only if
\begin{equation}\label{eq:frobenius}
u_{[a}\nabla_b u_{c]} = 0\,.
\end{equation}
\end{theorem}

\noindent
The connection to vorticity is direct: expanding
Eq.~\eqref{eq:frobenius} using the kinematic decomposition
shows that it is equivalent to $\omega_{ab} = 0$.

When orthogonal hypersurfaces exist, the induced three-metric
on $\Sigma_t$ is $\gamma_{ij} = h_{ab}\,e^a_i\,e^b_j$
(where $e^a_i$ are basis vectors tangent to $\Sigma_t$), and
we can define the \emph{three-dimensional Riemann tensor}
${}^{(3)}\!R_{ijkl}$ from the Levi-Civita connection of
$\gamma_{ij}$.  The three-dimensional Ricci scalar is
\begin{equation}\label{eq:3Ricci}
\sR := \gamma^{ij}\,{}^{(3)}\!R_{ij}\,,
\end{equation}
which encodes the intrinsic curvature of the spatial sections.
For an FLRW spacetime with curvature parameter $k$ and scale
factor $a$, we have $\sR = 6k/a^2$.

The relation between the four-dimensional and
three-dimensional curvatures is given by the Gauss equation:
\begin{equation}\label{eq:gauss}
{}^{(3)}\!R_{ijkl}
= h^a{}_i\,h^b{}_j\,h^c{}_k\,h^d{}_l\,R_{abcd}
  + K_{ik}\,K_{jl} - K_{il}\,K_{jk}\,,
\end{equation}
where $K_{ij} = h^a{}_i\,h^b{}_j\,\nabla_a u_b$ is the
extrinsic curvature of $\Sigma_t$.  In the irrotational case,
$K_{ij} = \frac{1}{3}\Theta\,\gamma_{ij} + \sigma_{ij}$
(the acceleration term drops from the spatial projection).


\subsection{The vortical case}\label{sec:vort-spatial}

When $\omega_{ab} \neq 0$, the Frobenius condition fails and
no family of orthogonal hypersurfaces exists.  Two approaches
are available.  In the \emph{threading} approach, we retain
the congruence $u^a$ and define ``spatial'' quantities by
projection with $h_{ab}$, even though the projected tensors do
not live on any genuine three-surface.  The ``spatial Ricci
scalar'' $\sR$ is then defined algebraically through the
constraint equations rather than through an intrinsic
three-geometry.  The threading approach is adopted by Ellis \& van
Elst~\cite{EllisVanElst1999} and the one we adopt.

In the \emph{slicing} approach, available when a Bianchi
symmetry group acts on the spacetime, the group orbits
provide natural spacelike hypersurfaces even when the fluid
congruence is not orthogonal to them.  The geometry-frame
congruence $u^a$ is chosen as the unit normal to these orbits,
and $\omega_{ab}^{(G)} = 0$ by construction.  Any non-zero
vorticity then appears only in the matter frame after the tilt
boost (Section~\ref{sec:frame-problem}).  For the Bianchi
framework, we therefore compute $\sR$ from the group-orbit
metric, where orthogonal surfaces do exist, and treat
matter-frame vorticity through the King--Ellis
transformation.

These two readings are kept separate as a machine-checkable contract: the
hypersurface-normal slicing carries $\Wstd = 0$ and admits the intrinsic
group-orbit Ricci substitution, whereas a threading/matter-frame vorticity term
is a \emph{blocked current result} until the full boost, energy-flux,
anisotropic-stress, and momentum-constraint terms are bound.  The split is
enforced by \texttt{FrameIdentityScope} (normal-frame versus threading
candidate); the corresponding Bianchi~VII$_h$ caveat is recorded in
Chapter~\ref{ch:bianchi-bounds} and the frame-vorticity repair note.


\subsection{Bianchi structure constants and $\sR$}
\label{sec:bianchi-Rstar}

For a spatially homogeneous spacetime, the Lie algebra of the
three-dimensional isometry group is specified by structure
constants $C^\alpha{}_{\beta\gamma}$.  Following the
Ellis--MacCallum~\cite{EllisMacCallum1969} parametrisation,
we decompose
\begin{equation}\label{eq:struct-decomp}
C^{\alpha}{}_{\beta\gamma}
= \epsilon_{\beta\gamma\delta}\,n^{\alpha\delta}
  + \delta^{\alpha}_{[\beta}\,a_{\gamma]}\,,
\end{equation}
where $n^{\alpha\beta} = n^{(\alpha\beta)}$ is a symmetric
$3\times 3$ tensor and $a_\alpha$ is a covector.  The Jacobi
identity requires $n^{\alpha\beta}a_\beta = 0$.  The
\emph{Class~A} models have $a_\alpha = 0$ (types I, II,
VI$_0$, VII$_0$, VIII, IX); the \emph{Class~B} models have
$a_\alpha \neq 0$ (types III, IV, V, VI$_h$, VII$_h$).

The spatial Ricci scalar for the group-orbit metric takes the
form~\cite{EllisMacCallum1969,RyanShepley1975}
\begin{equation}\label{eq:Rstar-bianchi}
\sR = -\tfrac{1}{2}\bigl(
n_\alpha{}^\alpha\, n_\beta{}^\beta
- n_{\alpha\beta}\,n^{\alpha\beta}\bigr)
+ 4\,a_\alpha\, a^\alpha\,.
\end{equation}
For Bianchi~I ($n^{\alpha\beta} = 0$, $a_\alpha = 0$), we have
$\sR = 0$: flat spatial sections.  For Bianchi~V
($n^{\alpha\beta} = 0$, $a_\alpha \neq 0$), we get
$\sR = 4\,a_\alpha a^\alpha > 0$ in the Ellis--MacCallum sign
convention, corresponding to $\Omega_k < 0$ (open, $k = -1$
FLRW curvature).  For Bianchi~IX
($n^{\alpha\beta} = \mathrm{diag}(1,1,1)$, $a_\alpha = 0$),
$\sR = -3 < 0$ in the normalised case, corresponding to
$\Omega_k > 0$ (closed, $k = +1$).


% ============================================================
\section{The generalised Friedmann equation}
\label{sec:gen-friedmann}
% ============================================================

Our goal in this section is to derive the Hamiltonian constraint
for a spatially homogeneous cosmology with vorticity---the
generalised Friedmann equation that replaces
Eq.~\eqref{eq:Friedmann1} of Chapter~\ref{ch:intro} when
anisotropy is present.

We start from the Gauss equation~\eqref{eq:gauss}.  Taking the
double trace $\gamma^{ik}\gamma^{jl}$ of the Gauss equation and
using the decomposition of $K_{ij}$ gives
\begin{equation}\label{eq:gauss-trace}
\sR = R + 2\,R_{ab}\,u^a u^b
  - K_{ij}K^{ij} + (K^i{}_i)^2\,.
\end{equation}
Now $K^i{}_i = \Theta$ and
$K_{ij}K^{ij} = \frac{1}{3}\Theta^2 + \sigma_{ab}\sigma^{ab}$
(since the cross term between the trace and the trace-free
part vanishes).  The Einstein equation gives
$R_{ab}\,u^a u^b = \frac{1}{2}\kappa(\rho + 3p) - \Lambda$
and $R = -\kappa(\rho - p) + 4\Lambda$.  Substituting and
simplifying:
\begin{equation}
\sR = -\kappa(\rho - p) + 4\Lambda
  + \kappa(\rho + 3p) - 2\Lambda
  - \tfrac{1}{3}\Theta^2 - 2|\sigma|^2 + \Theta^2\,.
\end{equation}
The $\rho$ and $p$ terms combine to $2\kappa\rho$, and
$-\frac{1}{3}\Theta^2 + \Theta^2 = \frac{2}{3}\Theta^2$.
Rearranging:
\begin{equation}\label{eq:ham-constraint-irrot}
\tfrac{1}{3}\Theta^2
= \kappa\,\rho + \Lambda
  - \tfrac{1}{2}\sR
  + |\sigma|^2\,.
\end{equation}
Equation~\eqref{eq:ham-constraint-irrot} is the Hamiltonian constraint for the irrotational case.
When vorticity is present, the threading approach introduces a
correction: the ``spatial Ricci scalar'' defined through the
constraint equations acquires a vorticity contribution, and the
generalised form becomes~\cite{EllisVanElst1999,EllisMaartens2012}
\begin{equation}\label{eq:ham-constraint}
\boxed{
\tfrac{1}{3}\Theta^2
= \kappa\,\rho_{\mathrm{tot}} + \Lambda
  - \tfrac{1}{2}\sR
  + |\sigma|^2 - |\omega|^2
}\,.
\end{equation}
The $-|\omega|^2$ term enters with the opposite sign to the
shear.  Vorticity acts as a \emph{negative} effective energy
density in the Friedmann equation.  For the geometry-frame
congruence in a Bianchi model, $\rho_{\mathrm{tot}}$ includes
the tilt contribution~\eqref{eq:rho-geom}.

Equation~\eqref{eq:ham-constraint} is the master constraint
from which the defect identity will be extracted.  It relates
the expansion rate to the total energy content, the spatial
curvature, and all kinematic degrees of freedom.


% ============================================================
\section{Standardised kinematic variables}\label{sec:std-vars}
% ============================================================

To compare kinematic quantities across models and epochs, we
normalise each scalar invariant by the expansion rate.  The
normalisation convention requires care, and we explain its
rationale.

\begin{definition}[Standardised variables]\label{def:std-vars}
The standardised shear, vorticity, and acceleration are
\begin{align}
\Sigstd &:= \frac{\sigma_{ab}\sigma^{ab}}{6\Htheta^2}\,,
\label{eq:Sigstd-def}\\
\Wstd   &:= \frac{\omega_{ab}\omega^{ab}}{6\Htheta^2}\,,
\label{eq:Wstd-def}\\
\Astd   &:= \frac{\dot{u}_a\dot{u}^a}{6\Htheta^2}\,,
\label{eq:Astd-def}
\end{align}
where $\Htheta := \Theta/3$ is the Hubble rate of the
geometry-frame congruence.  Each variable is dimensionless,
non-negative, and vanishes if and only if the corresponding
kinematic quantity vanishes.
\end{definition}

The factor of six in the denominator is a convention choice.
We adopt it so that the conversion between the standardised
variable and the shear-to-expansion ratio takes the simple form
$\sigma_{ab}\sigma^{ab}/\Theta^2 = (2/3)\Sigstd$, or
equivalently $(\sigma/\Theta)^2 = (2/3)\Sigstd$ where
$\sigma/\Theta := \sqrt{\sigma_{ab}\sigma^{ab}}/\Theta$.
The MES shear bound~\eqref{eq:intro-MES} constrains
$\sigma/\Theta < \Bsig$, which translates to
$\Sigstd < \tfrac{3}{2}\Bsig^2$.

For vorticity, two conventions coexist in the literature:
$\omega/H := \sqrt{\omega_a\omega^a}/H$
(Saadeh~et~al.~\cite{Saadeh2016b}) and
$\bar{\omega} := \sqrt{\omega_{ab}\omega^{ab}}/\Theta$.
The relations to $\Wstd$ are
\begin{equation}\label{eq:omH-W2}
(\omega/H)^2 = 3\,\Wstd\,,\qquad
\bar{\omega}^2 = \tfrac{2}{3}\,\Wstd\,.
\end{equation}
The Saadeh~et~al.\ 95\% upper limit
$(\omega/H)_0 < 5.2 \times 10^{-11}$ thus translates to
$\Wstd < 9.01 \times 10^{-22}$.


\subsection{Density parameters}\label{sec:density-params}

Dividing the generalised Friedmann
equation~\eqref{eq:ham-constraint} by $3\Htheta^2$ converts
every term into a dimensionless density parameter.  We define
\begin{equation}\label{eq:density-params}
\Omega_{\mathrm{tot}} := \frac{\kappa\,\rho_{\mathrm{tot}}}{3\Htheta^2}\,,\qquad
\OL := \frac{\Lambda}{3\Htheta^2}\,,\qquad
\Omega_k := -\frac{\sR}{6\Htheta^2}\,,
\end{equation}
so that the closure relation reads
\begin{equation}\label{eq:closure-full}
1 = \Omega_{\mathrm{tot}} + \OL + \Omega_k
  + \Sigstd - \Wstd\,.
\end{equation}
Equation~\eqref{eq:closure-full} is the normalised Hamiltonian constraint.  Every Bianchi
model, at every instant, satisfies this relation exactly.


% ============================================================
\section{FLRW comparators and the departure parameter}
\label{sec:defect-def}
% ============================================================

\subsection{Comparator FLRW models}\label{sec:comparators}

To define the ``distance from FLRW'' of a Bianchi model, we
need a reference FLRW geometry against which to measure the
departure.  We call this the \emph{comparator}.  A comparator
is specified by choosing a reference curvature parameter
$\Okref$ and defining the FLRW matter budget as
\begin{equation}\label{eq:FLRW-budget}
\OFLRW := 1 - \OL - \Okref\,.
\end{equation}
Three comparator policies are natural:
\begin{equation}\label{eq:comparator-policies}
\compflat:\;\; \Okref = 0\,,\quad
\compclosed:\;\; \Okref = \Omega_k^{(\mathrm{IX})}\,,\quad
\compmatched:\;\; \Okref = \Omega_k\,.
\end{equation}
The \emph{flat comparator} $\compflat$ ($\Okref = 0$) is
appropriate for types with vanishing spatial curvature (Bianchi~I
and its tilted extensions).  The \emph{closed comparator}
$\compclosed$ uses the Bianchi~IX isotropic curvature.  The
\emph{curvature-matched comparator} $\compmatched$
($\Okref = \Omega_k$) sets the anisotropic curvature
contribution to zero and isolates the shear and tilt
contributions.  Throughout this work, we use $\compflat$ unless
otherwise stated.

\begin{center}
\fbox{\parbox{0.85\textwidth}{\small
\textbf{Warning.}  The three comparators $\compflat$,
$\compclosed$, $\compmatched$ are \emph{not} identical: they
assign different values of $x$ to the same spacetime.  In
particular, $x = 0$ under $\compflat$ means the model coincides
with a spatially flat FLRW, while $x = 0$ under $\compmatched$
means only that the shear and tilt vanish---the model may still
have non-zero spatial curvature.  \emph{Every} statement
involving $x$, $Q$, $\Pi$, or $\FC$ is implicitly conditioned on
the comparator.}}
\end{center}


\subsection{The departure parameter}\label{sec:defect-x}

\begin{definition}[Departure parameter]\label{def:defect}
The \emph{defect} of a Bianchi model relative to a specified
FLRW comparator is
\begin{equation}\label{eq:defect-def}
x := \OFLRW - \Omega\,,
\end{equation}
where $\Omega := \kappa\hat{\rho}/(3\Htheta^2)$ is the
rest-frame (matter-frame) density parameter and $\hat{\rho}$
is the matter-frame energy density.
\end{definition}

The departure parameter $x$ measures how much the matter density falls short
of the FLRW budget.  Every unit of shear, tilt, or anisotropic
curvature diverts energy density away from the FLRW value,
increasing $x$.

\begin{proposition}[Sector-dependent sign structure of $x$ {\normalfont $[\ExactStatus]$}]
\label{prop:x-sign}
The departure parameter $x$ is a signed quantity.  Its sign is determined by
the sector to which the model belongs.
\begin{enumerate}
\item[\emph{(i)}]
  \textbf{Irrotational non-negative sector} ($\Wstd = 0$,\;
  $\Okaniso \geq 0$):
  \begin{equation}\label{eq:x-irrot-pos}
  x = \Sigstd + \Otilt + \Okaniso \geq 0\,,
  \end{equation}
  with equality if and only if the model coincides with the
  chosen FLRW comparator.  All flat orthogonal Bianchi types
  (BI, BVII$_0$, BII, BVI$_0$) and all flat tilted types (BI,
  BIII) under the flat comparator belong to this sector.
  Certification of $\FC = x/\xmax$ is straightforward.

\item[\emph{(ii)}]
  \textbf{Irrotational negative sector} ($\Wstd = 0$,\;
  $\Okaniso < 0$ dominant):
  \begin{equation}\label{eq:x-irrot-neg}
  x = \Sigstd + \Otilt + \Okaniso < 0
  \quad\text{when}\quad
  |\Okaniso| > \Sigstd + \Otilt\,.
  \end{equation}
  The Bianchi~IX tilted model is the principal example:
  $x = -1.36 \times 10^{-3}$, driven entirely by
  $\Okaniso = -1.36 \times 10^{-3}$ from the Planck curvature
  constraint (Table~\ref{tab:departure_summary}).  The standard
  filling fraction $\FC = x/\xmax$ is ill-defined; a signed
  saturation coordinate is required.  BVIII and BIX~(orth) under
  the flat comparator also belong to this sector.

\item[\emph{(iii)}]
  \textbf{Vortical sector} ($\Wstd > 0$):
  \begin{equation}\label{eq:x-vort}
  x = \Sigstd - \Wstd + \Otilt + \Okaniso\,,
  \end{equation}
  where the sign depends on the balance of all four terms.
  Certification is blocked: the $\Wstd$ subtraction can make $x$
  negative even when the geometry is far from FLRW.  The
  BVII$_h$ tilted models belong to this sector.
\end{enumerate}
\end{proposition}

\noindent
The proof of part~(i) is immediate from
Theorem~\ref{thm:master-defect-id}: every term is individually
non-negative.  Part~(ii) follows from the dominance of the
negative $\Okaniso$ term.  Part~(iii) follows from the
negative sign of $\Wstd$ in the identity.  The three sectors
are exhaustive and mutually exclusive; every Bianchi model
belongs to exactly one sector under a given comparator.


% ============================================================
\section{The master defect identity}\label{sec:master-id}
% ============================================================

We now state and prove the central algebraic result of this
framework.

\begin{theorem}[Extended master defect identity {\normalfont $[\ExactStatus]$}]
\label{thm:master-defect-id}
Let $(\mathcal{M}, g_{ab})$ be a spatially homogeneous
(Bianchi) cosmology with geometry-frame congruence $u^a$
[$\Gframe$-frame],
expansion $\Theta > 0$, kinematical Hubble rate
$\Htheta = \Theta/3$, shear $\sigma_{ab}$ [$\Gframe$], and
vorticity $\omega_{ab}$ [$\Gframe$].  Assume:
\begin{enumerate}
\item[\emph{(i)}] $u^a$ is the Bianchi group-invariant time
  direction and $\sR$ is computed from the group-orbit slicing.
\item[\emph{(ii)}] The generalised Hamiltonian constraint
  \eqref{eq:ham-constraint} holds.
\item[\emph{(iii)}] An FLRW comparator is specified by
  $\OL$ and $\Okref$, with
  $1 = \OFLRW + \OL + \Okref$.
\item[\emph{(iv)}] The matter content is a perfect fluid with
  equation of state $w > -1$ and tilt rapidity $\beta$.
\end{enumerate}
Then the departure parameter $x := \OFLRW - \Omega$ satisfies
\begin{equation}\label{eq:master-identity}
\boxed{x = \Sigstd - \Wstd + \Otilt + \Okaniso}\,,
\end{equation}
where each term is defined as follows:
\begin{align}
\Sigstd &= \frac{\sigma_{ab}\sigma^{ab}}{6\Htheta^2} \geq 0
  &&\text{(standardised shear)}\,,
  \label{eq:Sigstd-id}\\[3pt]
\Wstd &= \frac{\omega_{ab}\omega^{ab}}{6\Htheta^2} \geq 0
  &&\text{(standardised vorticity)}\,,
  \label{eq:Wstd-id}\\[3pt]
\Otilt &= (1+w)\,\Omega\,\sinh^2\!\beta \geq 0
  &&\text{(tilt density)}\,,
  \label{eq:Otilt-def}\\[3pt]
\Okaniso &= \Omega_k - \Okref
  &&\text{(anisotropic curvature)}\,.
  \label{eq:Oka-def}
\end{align}
\end{theorem}

\begin{proof}
The argument proceeds in four steps from the normalised
Hamiltonian constraint.

\emph{Step~1 (Tilt split).}
The geometry-frame total energy density decomposes via
Eq.~\eqref{eq:rho-geom} as
$\rho_{\mathrm{tot}} = \hat{\rho}
+ (\hat{\rho} + \hat{p})\sinh^2\!\beta$.
Dividing by $3\Htheta^2$ and using
$\hat{p} = w\hat{\rho}$, we obtain
\begin{equation}\label{eq:Otot-split}
\Omega_{\mathrm{tot}}
= \Omega + (1+w)\,\Omega\,\sinh^2\!\beta
= \Omega + \Otilt\,,
\end{equation}
where $\Omega = \kappa\hat{\rho}/(3\Htheta^2)$ is the
matter-frame density parameter.  The condition $w > -1$ and
$\Omega \geq 0$ ensures $\Otilt \geq 0$.

\emph{Step~2 (Substitution into closure).}
We substitute Eq.~\eqref{eq:Otot-split} into the closure
relation~\eqref{eq:closure-full}:
\begin{equation}\label{eq:pf-step2}
1 = \Omega + \Otilt + \OL + \Omega_k
  + \Sigstd - \Wstd\,.
\end{equation}

\emph{Step~3 (Curvature split).}
We write $\Omega_k = \Okref + \Okaniso$ by
Definition~\eqref{eq:Oka-def}, so that
\begin{equation}\label{eq:pf-step3}
1 = \Omega + \Otilt + \OL + \Okref + \Okaniso
  + \Sigstd - \Wstd\,.
\end{equation}

\emph{Step~4 (Comparator subtraction).}
The comparator FLRW satisfies $1 = \OFLRW + \OL + \Okref$
by Eq.~\eqref{eq:FLRW-budget}.  Subtracting this from
Eq.~\eqref{eq:pf-step3} and rearranging:
\begin{equation}
\OFLRW - \Omega
= \Sigstd - \Wstd + \Otilt + \Okaniso\,.
\end{equation}
Since $x = \OFLRW - \Omega$ by
Definition~\ref{def:defect}, the identity follows.
\end{proof}

We pause to emphasise several properties of this result.
the conditioned legacy figure gallery in Appendix~\ref{app:conditioned-legacy-figure-gallery} provides a schematic
decomposition of the four terms.

\begin{remark}[Identity, not inequality]\label{rem:identity}
Equation~\eqref{eq:master-identity} is exact and algebraic:
it holds at every instant without perturbative or near-FLRW
approximation.  Inequalities arise only when individual terms
are bounded using observational data.
\end{remark}

\begin{remark}[FLRW limit]\label{rem:FLRW-limit}
We verify the FLRW limit explicitly.  An FLRW model has
$\sigma_{ab} = 0$ (hence $\Sigstd = 0$),
$\omega_{ab} = 0$ (hence $\Wstd = 0$),
$\beta = 0$ (hence $\Otilt = 0$), and isotropic curvature.
If the comparator is curvature-matched ($\Okref = \Omega_k$),
then $\Okaniso = 0$ and the identity gives
$x = 0 + 0 + 0 + 0 = 0$.
For the flat comparator, an FLRW model with $\Omega_k \neq 0$
gives $x = \Okaniso \neq 0$; thus $x = 0$ with a flat
comparator implies the model is \emph{spatially flat} FLRW.
In short, $x = 0$ means the model coincides with the chosen
comparator---not with ``FLRW'' in an absolute sense.
\end{remark}

\begin{remark}[Sign of the vorticity term]
\label{rem:vort-sign}
Vorticity enters with a \emph{negative} sign, reflecting its
role in the generalised Friedmann equation as a negative
effective energy density.  For observationally relevant
amplitudes ($\Wstd < 10^{-21}$ from
Eq.~\eqref{eq:omH-W2} and the Saadeh bound), this term is
smaller than $\Sigstd \sim 10^{-6}$ by fifteen orders of
magnitude; its contribution to $x$ is negligible in all
irrotational and weakly-vortical models.

The sign has stronger consequences when combined with negative
anisotropic curvature.  For the Bianchi~IX tilted model, the
dominant contribution is $\Okaniso = -1.36 \times 10^{-3}$
(from the Planck curvature constraint), giving $x = -1.36
\times 10^{-3}$---an explicitly negative defect
(Proposition~\ref{prop:x-sign}, part~ii).  The occupancy ratio
$Q = x / x_{\max}$ inherits this sign, producing $Q = -211$
and exceedance $\Pi(q_\star) = 0$ for all positive thresholds.
This is not a failure of the framework: negative $x$ means the
model \emph{overshoots} the FLRW comparator, and the
three-layer structure correctly propagates this through $Q$ and
$\Pi$ without artificial sign-clipping.
\end{remark}

\begin{remark}[Irrotational limit]\label{rem:irrot}
Setting $\omega_{ab} = 0$ gives $\Wstd = 0$ and reduces the
identity to $x = \Sigstd + \Otilt + \Okaniso$---a sum of three
non-negative terms.  In this limit, the component-wise bounds
$x \geq \Sigstd$, $x \geq \Otilt$, and $x \geq \Okaniso$
follow immediately.
\end{remark}

\begin{corollary}[Component-wise bounds {\normalfont $[\ConditionalStatus]$}]
\label{cor:component-bounds}
In the irrotational sector ($\Wstd = 0$),
$x \geq \Sigstd$, $x \geq \Otilt$, and
$x \geq \Okaniso$, with equality if and only if the remaining
two terms vanish.
\end{corollary}

\begin{corollary}[Shear--expansion conversion]
\label{cor:sigma-from-Bsigma}
The MES bound $\sigma/\Theta < \Bsig$ implies
$\Sigstd < \frac{3}{2}\Bsig^2$.
\end{corollary}
\begin{proof}
By definition, $(\sigma/\Theta)^2
= \sigma_{ab}\sigma^{ab}/\Theta^2 = (2/3)\Sigstd$.
The bound $\sigma/\Theta < \Bsig$ gives
$(2/3)\Sigstd < \Bsig^2$, hence
$\Sigstd < \frac{3}{2}\Bsig^2$.
\end{proof}

\begin{remark}[Frame attribution of the defect identity]
\label{rem:frame-ledger}
The master identity~\eqref{eq:master-identity} mixes
quantities from different frames.
Table~\ref{tab:frame-ledger} records the frame
origin of each term.

\begin{table}[ht]
\caption{Frame ledger for the master defect identity.
\Gframe\ = geometry frame; \Mframe\ = matter frame;
\GMframe\ = mixed (involves M$\to$G boost).}
\label{tab:frame-ledger}
\centering\footnotesize
\begin{tabular}{lcl}
\hline\hline
Quantity & Frame & Origin \\
\hline
$\Sigstd$ & \Gframe & $\sigma_{ab}\sigma^{ab}/(6\Htheta^2)$ \\
$\Wstd$ & \Gframe & $\omega_{ab}\omega^{ab}/(6\Htheta^2)$ \\
$\Okaniso$ & \Gframe & $\Omega_k - \Okref$ \\
$\Otilt$ & \GMframe & $(1{+}w)\Omega\sinh^2\!\beta$
  (involves M$\to$G rapidity) \\
$x$ & \GMframe & $\Sigstd - \Wstd + \Otilt + \Okaniso$
  (hybrid via $\Otilt$) \\
\hline\hline
\end{tabular}
\end{table}

\noindent
The defect $x$ is \emph{not} a pure geometry-frame scalar: it
inherits mixed-frame character from $\Otilt$, which depends on
the tilt rapidity $\beta$ between the matter and geometry
congruences.  For irrotational, orthogonal models
($\beta = 0$, $\omega_{ab} = 0$), the identity reduces to
$x = \Sigstd + \Okaniso$, which is pure \Gframe.  For tilted
models, the \GMframe\ character is irreducible.
\end{remark}


% ============================================================
\section{The three-layer diagnostic pushforward framework}
\label{sec:three-layer}
% ============================================================

Our goal in this section is to connect the exact algebraic
identity~\eqref{eq:master-identity} to a registered diagnostic
pushforward.  In HTT-owned likelihood contexts the same objects
may be evaluated on posterior samples, but MIO reports only
diagnostic certificates and registered curves.  The departure
parameter $x$ is not directly observable: it is a function of
the model parameters $\theta$ through the master identity.  What
enters the statistical analysis is the typed pushforward
distribution of $x$ induced by propagating the declared input
measure through the identity.
We organise this into three layers, each providing a distinct
level of interpretive resolution.


\subsection{Layer~1: The pushforward defect}
\label{sec:layer1}

Given posterior samples $\{\theta_i\}_{i=1}^N$ from the nested
sampling of any Bianchi model, we compute
\begin{equation}\label{eq:xi-pushforward}
x_i = \Sigstd(\theta_i) - \Wstd(\theta_i)
      + \Otilt(\theta_i) + \Okaniso(\theta_i)
\end{equation}
for each sample, inheriting the posterior weights $w_i$.  The
resulting weighted collection $\{x_i, w_i\}$ is the declared
pushforward sample for $x$---exact with respect to the stated
input samples, since the identity holds for every $\theta$.

The component-wise decomposition is retained: each sample
carries all four contributions, enabling statements such as
``the posterior median tilt contribution is
$\tilde{\Otilt} = 5.85 \times 10^{-7}$, while the posterior
median shear contribution is
$\tilde{\Sigstd} \lesssim 10^{-18}$''
(FLRW$_{\mathrm{tilt}}$ model, flat comparator, VER05 corrected; see
Table~\ref{tab:departure_summary}).  The shear is negligible by
nineteen orders of magnitude: the entire defect is tilt-driven.

\begin{definition}[Pushforward diagnostic summaries]
\label{def:x-summaries}
We report:
the posterior median $\tilde{x}$,
the posterior mean $\bar{x}$,
the 95\% highest posterior density interval
$[x_{2.5}, x_{97.5}]$,
and the component medians
$\tilde{\Sigstd}$, $\tilde{\Wstd}$,
$\tilde{\Otilt}$, $\tilde{\Okaniso}$.
\end{definition}


\subsection{Layer~2: The occupancy ratio}
\label{sec:layer2}

The defect $x$ depends on the normalisation of the Hubble rate.
To obtain a dimensionless measure of ``how close to the
ceiling,'' we define the \emph{occupancy ratio}
\begin{equation}\label{eq:Q-def}
Q := \frac{x}{x_{\max}}\,,
\end{equation}
where $x_{\max}$ is the MES algebraic ceiling---the largest
value of $\Sigstd$ permitted by the MES shear bound for
the adopted scenario.
For the S3 kinematic scenario with $\epsilon_1 = 1.233
\times 10^{-3}$, the ceiling is
$x_{\max} = \frac{3}{2}\Bsig^2 = 6.47 \times 10^{-6}$
(Chapter~\ref{ch:bianchi-bounds},
Table~\ref{tab:bounds}).

Because the numerator $x$ is a departure (tilt-sector dominated)
projection while the denominator $x_{\max}$ is a shear-sector
ceiling, the scalar $Q$ is a \emph{channel-mismatched proxy score
only}, not a channel-matched physical occupancy.  Channel-matched
occupancy is reserved for the MIO channel-matched occupancy vector,
whose numerator and denominator share a sector; the classifier
\texttt{classify\_occupancy\_language} blocks ``physical occupancy''
wording for scalar $Q/F$ rows unless the channels match or a joint
admissible-ceiling proof is attached (see
\texttt{docs/generated/channel\_occupancy\_repair\_report.md}).

The HTT posterior use of $Q$ is a ratio of two model-dependent
quantities evaluated under a declared transfer path, comparator,
and denominator policy.  MIO may report the same algebraic ratio
as a diagnostic coordinate, but the MIO diagnostic $Q$ has no
posterior, occupancy, or evidence semantics.  In both lanes
$Q$ inherits the sign of $x$: for models with
$x < 0$ (such as Bianchi~IX with dominant negative $\Okaniso$),
$Q < 0$.

For the FLRW$_{\mathrm{tilt}}$ model under the flat comparator,
the posterior mean occupancy is $\bar{Q} = 0.092$, with a 95\%
highest posterior density interval
$[0.035,\, 0.19]$.  That is, the data place the defect at
roughly $9\%$ of the algebraic ceiling in this conditioned HTT
posterior pushforward---resolved within the legacy diagnostic,
but far from saturation.


\subsection{Layer~3: The HTT posterior exceedance}
\label{sec:layer3}

The occupancy ratio $Q$ records how a signed projection compares
with a declared denominator policy.  To make threshold use
auditable, we define the registered HTT \emph{posterior
exceedance curve}
\begin{equation}\label{eq:Pi-def}
\Pi_{\rm HTT}(q_\star) := P(Q > q_\star \,|\, D)\,,
\end{equation}
where $q_\star$ is a chosen threshold.  For the
FLRW$_{\mathrm{tilt}}$ model: $\Pi_{\rm HTT}(0.01) = 1.000$
(the posterior unambiguously exceeds $1\%$ of the ceiling),
$\Pi_{\rm HTT}(0.05) = 0.974$,
$\Pi_{\rm HTT}(0.1) = 0.352$, and
$\Pi_{\rm HTT}(0.5) = 0.000$.  This legacy conditioned example exceeds
the $1$--$5\%$ thresholds under its recorded inputs, is
indeterminate at $10\%$, and does not exceed $50\%$.  The curve
is not a truth probability, p-value, or geometry-detection
probability.
It is an HTT posterior exceedance probability for this legacy
conditioned run, distinct from the MIO diagnostic exceedance curve
$\Pi_{\rm MIO}$, which is an empirical exceedance fraction with
no posterior, occupancy, or evidence semantics.

\begin{remark}[Interpretation guide]\label{rem:Pi-interpretation}
The three layers form a chain of increasing interpretive
commitment.  Layer~1 ($x$) is exact algebra---no choice is made
beyond the comparator.  Layer~2 ($Q$) introduces the MES
ceiling, coupling the result to the adopted observational bound.
Layer~3 ($\Pi$) introduces a threshold $q_\star$, whose choice
must be registered before interpretation and accompanied by
look-elsewhere and null-status metadata.
A referee who objects to the ceiling choice can read $x$ alone;
a referee who accepts the ceiling but disputes the threshold
can read $Q$ alone.  This stratification makes the
interpretive dependencies explicit and auditable.
\end{remark}


\subsection{Comparator sensitivity of the three layers}
\label{sec:comp-sensitivity}

We verified numerically that, for all models without a curvature
parameter ($\Okaniso = 0$ by construction), the three layers are
identically invariant under the comparator choice: the flat,
closed, and curvature-matched comparators yield the same
$x$, $Q$, and $\Pi$ to machine precision.  This follows from
the fact that $\Okaniso$ is the only comparator-dependent term
in the identity.

For models with a curvature parameter (Bianchi~V and~IX),
comparator dependence is maximal: the exceedance spread is
$\Delta_{\Pi}^{\mathrm{comp}}(0.01) = 1.0$ for both
BV$_{\mathrm{tilt}}$ and BIX$_{\mathrm{tilt}}$
(Table~\ref{tab:departure_summary}).  This is expected:
the curvature contribution $\Okaniso$ switches sign between
comparators, changing the qualitative conclusion.  We report
all departure quantities under the flat comparator throughout
and flag curvature-bearing models explicitly.


\subsection{Ontological status of derived quantities}
\label{sec:ontology}

The three-layer framework reports two additional derived
quantities that must be carefully distinguished from the
departure components.

The \emph{tilt velocity}
$v_{\mathrm{tilt}} := c\,\tanh\beta$ is a physical
observable---the peculiar velocity of the matter frame
relative to the geometry frame.  For the FLRW$_{\mathrm{tilt}}$
posterior,
$\tilde{v}_{\mathrm{tilt}} = 408\;\mathrm{km\,s}^{-1}$.
This quantity is directly comparable to the CF4 bulk-flow
measurement.

The \emph{tilt density} $\Otilt = (1+w)\Omega\sinh^2\!\beta$
is a component of the departure identity---a bookkeeping
variable that enters the algebraic constraint, not a directly
measurable energy density.  Its value
($\tilde{\Otilt} = 5.85 \times 10^{-7}$) should not be
interpreted as a ``fifth energy component'' but rather as the
portion of the FLRW matter budget redirected by the tilt into
the geometry-frame total.

We enforce this distinction throughout: $v_{\mathrm{tilt}}$
appears in observational comparisons
(Chapter~\ref{ch:discussion}), while $\Otilt$ appears only in
the defect identity algebra and its decomposition figures.


% ============================================================
\subsection{Ceiling semantics and certified \texorpdfstring{$\FC$}{F\_C}}
\label{sec:ceiling-semantics}
% ============================================================

The three-layer framework uses $\xmax$ as a normalisation
ceiling for the occupancy ratio $Q = x/\xmax$.  We now formalise
the conditions under which this normalisation carries physical
content.

\begin{definition}[Admissible ceiling family]
\label{def:ceiling-family}
An \emph{admissible ceiling family} is a map
$\xmaxC: \Theta \to \mathbb{R}_{>0}$,
where $\Theta$ is the parameter space and $C$ labels the
comparator, such that $\xmaxC(\theta) \geq |x(\theta)|$
for all $\theta$ in the sector where $\xmaxC$ is
evaluated.  The MES ceiling $\xmax = \frac{3}{2}\Bsig^2$
is the principal instance used in this work.
\end{definition}

\begin{definition}[Class-conditioned filling fraction]
\label{def:FC}
For a model in sign-domain sector $S$ under comparator $C$,
the \emph{class-conditioned filling fraction} is
\begin{equation}\label{eq:FC-def}
\FC := \frac{x}{\xmaxC}\,,
\end{equation}
defined only when $x \geq 0$ (irrotational non-negative
sector) and $\xmaxC > 0$.
\end{definition}

\begin{proposition}[No automatic physical interpretation {\normalfont $[\ExactStatus]$}]
\label{prop:FC-no-auto}
The class-conditioned filling fraction $\FC$ does \emph{not}
automatically carry a physical interpretation as ``the fraction
of the anisotropy budget that is filled.''  A counterexample is
the Bianchi~IX orthogonal model under the flat comparator: the
negative $\Okaniso = -1.36 \times 10^{-3}$ drives $x < 0$,
making $\FC$ undefined by Definition~\ref{def:FC}.  Switching
to $\compmatched$ sets $\Okaniso = 0$ and restores $\FC \geq 0$,
but the ``physical meaning'' of $\FC$ now depends entirely on
the comparator choice.
\end{proposition}

\begin{proposition}[Certified saturation in the sign-clean
  sector {\normalfont $[\ConditionalStatus]$}]\label{prop:FC-certified}
In the irrotational non-negative sector
(Proposition~\ref{prop:x-sign}, part~(i)), under any admissible
ceiling, $0 \leq \FC \leq 1$.  The value $\FC = 0$ means exact
coincidence with the comparator; $\FC = 1$ means the defect
saturates the ceiling.  The MES ceiling gives $\FC \in
[0, 0.090]$ (WFH2009 primary) to $[0, 0.063]$ (CF4++ sensitivity variant)
across the observationally relevant scenarios.
\end{proposition}

\begin{remark}[Signed saturation coordinate]
\label{rem:signed-sat}
For models in the irrotational negative or vortical sectors,
a signed saturation coordinate $\FC^{\pm} := x / |\xmaxC|$
can be defined, taking values on $(-\infty, 1]$.  Negative
values indicate ``overshoot'' beyond the comparator, not a
failure of the framework.  We report $\FC^{\pm}$ for these
models in the appendix but do not include them in the main
evidence or occupancy tables.
\end{remark}


% ── Sector partition and conditional diagnostics ─────────
\subsection{Sector partition and posterior decomposition}
\label{sec:sector-partition}

The sign structure of $x$ partitions the state space into
three sectors that govern which diagnostics are well-defined.

\begin{definition}[Sign-domain sector partition]
\label{def:sector-partition}
For each Bianchi type $T$ and comparator $C$, the constrained
state space $\Theta_T$ admits the measurable partition
\begin{equation}\label{eq:sector-partition}
\Theta_T
= \Theta_{T,+}^{\mathrm{irr}}
  \;\sqcup\;
  \Theta_{T,-}^{\mathrm{irr}}
  \;\sqcup\;
  \Theta_T^{\mathrm{vort}}\,,
\end{equation}
where
$\Theta_{T,+}^{\mathrm{irr}} := \{\theta : W^2 = 0,\;
x(\theta) \geq 0\}$ (irrotational non-negative),
$\Theta_{T,-}^{\mathrm{irr}} := \{\theta : W^2 = 0,\;
x(\theta) < 0\}$ (irrotational negative),
and $\Theta_T^{\mathrm{vort}} := \{\theta : W^2 > 0\}$
(vortical).
\end{definition}

\begin{definition}[Posterior sector masses]
\label{def:sector-mass}
Given the posterior $p(\theta \,|\, D, T)$, the
\emph{posterior sector mass} for sector $s$ is
\begin{equation}\label{eq:sector-mass}
m_{T,s}(D)
:= \int_{\Theta_{T,s}} p(\theta \,|\, D, T)\,
   d\theta\,.
\end{equation}
The three masses sum to unity.  A large $m_{T,-}^{\mathrm{irr}}$
signals that the posterior places substantial weight in the
negative-departure region, where the occupancy ratio $\FC$ is
undefined (Definition~\ref{def:FC}) and the signed coordinate
$\FC^{\pm}$ (Remark~\ref{rem:signed-sat}) must be used instead.
\end{definition}

The full posterior summary for a model is therefore not a single
occupancy number but the triple
\begin{equation}\label{eq:full-summary}
\bigl(\,
  m_{T,+}^{\mathrm{irr}}(D),\;
  m_{T,-}^{\mathrm{irr}}(D),\;
  m_T^{\mathrm{vort}}(D)\,\bigr)
\end{equation}
together with the sector-conditional diagnostic laws
$p(\FC \,|\, D,\,\theta \in \Theta_{T,+}^{\mathrm{irr}})$ and
$p(\FC^{\pm} \,|\, D,\,\theta \in \Theta_{T,-}^{\mathrm{irr}}
\cup \Theta_T^{\mathrm{vort}})$.

\begin{proposition}[Sectoral ceiling existence
  {\normalfont $[\ConditionalStatus]$}]
\label{prop:ceiling-existence}
Let $C_{\mathrm{irr}}^{+}$ denote the irrotational non-negative
sector under the flat comparator, and suppose:
\textup{(i)}~the MES shear bound gives
$\Sigstd \leq (3/2)\Bsig^2$ for the adopted $\eps_\ell$;
\textup{(ii)}~the safe-route tilt bound gives
$\beta \leq \eps_1/(1 + \etaud)$;
\textup{(iii)}~$\Omega_m \in (0,1]$ and $w \in [0, 1/3]$;
\textup{(iv)}~$C_{\mathrm{irr}}^{+}$ is a closed bounded
subset of $\Theta_T$.
Then
\begin{equation}
U_{C_{\mathrm{irr}}^{+}}
:= \max_{\theta \in C_{\mathrm{irr}}^{+}} x(\theta)
= \tfrac{3}{2}\Bsig^2 + (1+w)\Omega_m\sinh^2\!\beta_{\max}
\end{equation}
exists and is finite.
\end{proposition}

\begin{proof}
In $C_{\mathrm{irr}}^{+}$ we have $W^2 = 0$ and
$\Okaniso \geq 0$, so $x = \Sigstd + \Otilt + \Okaniso$.
Each summand is non-negative and bounded above:
$\Sigstd \leq (3/2)\Bsig^2$ by~(i),
$\Otilt = (1+w)\Omega_m\sinh^2\!\beta
\leq (1+w)\Omega_m\sinh^2\!\beta_{\max}$ by~(ii)--(iii),
and $\Okaniso \leq (3/2)\Bsig^2$ (it contributes to the
same MES budget).  Hence $x$ is continuous and bounded on a
compact set; by the Weierstrass extreme-value theorem the
maximum exists.
\end{proof}

\begin{remark}[Scope of the ceiling guarantee]
\label{rem:ceiling-scope}
Proposition~\ref{prop:ceiling-existence} does \emph{not}
extend to the irrotational negative or vortical sectors.  In
those sectors, the interplay of $-W^2$ and negative $\Okaniso$
can make $\sup x$ depend on cancellation structures for which
no universal bound exists.  For this reason the occupancy ratio
$Q = x/\xmax$ carries $\ConditionalStatus$: it is a
theorem-backed quantity only within $C_{\mathrm{irr}}^{+}$, and
a proxy score elsewhere.
\end{remark}


% ── Identifiability under scalar-amplitude compression ────
\subsection{Identifiability under scalar-amplitude compression}
\label{sec:identifiability}

The scalar-amplitude likelihood compresses all directional
information into amplitude summary statistics.  We formalise
which quantities survive this compression.

\begin{definition}[$K$-identifiable functional]
\label{def:K-ident}
Let $K_{M,\theta}(z)$ denote the likelihood evaluated on
compressed observables $z$ for model $M$ at parameters~$\theta$.
A functional $g(M,\theta)$ is \emph{$K$-identifiable} if
\begin{equation}\label{eq:K-ident}
K_{M,\theta}(z) = K_{M',\theta'}(z)\;\forall z
\quad\Longrightarrow\quad
g(M,\theta) = g(M',\theta')\,.
\end{equation}
Equivalently, $g$ factors through the quotient
$\mathcal{Q}_K := \Theta / {\sim_K}$.
\end{definition}

\begin{proposition}[$K$-identifiability audit]
\label{prop:K-audit}
Under the VER06 scalar-amplitude likelihood:

\textup{(i)}~$\beta$, $v_{\mathrm{tilt}} = c\tanh\beta$,
$\sinh^2\!\beta$, $\Otilt$ are $K$-identifiable
\emph{and data-constrained}: the posterior is data-dominated
with $N_{\mathrm{eff}} > 700$.

\textup{(ii)}~$\Sigstd$ is $K$-identifiable (constant within
each equivalence class) but \emph{not data-constrained}:
the posterior fills the prior at $\sim\!10^{-20}$.

\textup{(iii)}~$x_C$ under the flat comparator for
curvature-free types is $K$-identifiable and partially
data-constrained (via $\Otilt$).  For curvature-bearing types,
$x_C$ is \emph{not} $K$-identifiable: $\Omega_K$ is inactive
and comparator-dependent.

\textup{(iv)}~Any geometry norm $\mathcal{G}_u$ that depends on
the \emph{type label} (beyond the equivalence class) is not
$K$-identifiable.  In particular, the distinction between
Bianchi~I and Bianchi~VII$_0$ orthogonal, or between
FLRW$_{\mathrm{tilt}}$ and BI$_{\mathrm{tilt}}$, is not
accessible at the scalar-amplitude level.
\end{proposition}

\begin{proof}
Parts~(i)--(iii) follow by inspection of the
\texttt{predicted\_observables()} output: within each
equivalence class (Section~\ref{sec:equivalence} of
Chapter~\ref{ch:results}), the likelihood-active parameters
are identical.  Part~(iv) is immediate: the geometry norm
$\mathcal{G}_u$ distinguishes types that share the same
$K_{M,\theta}$, hence it does not factor through
$\mathcal{Q}_K$.
\end{proof}

\begin{remark}[Geometry claim prohibition as a theorem]
\label{rem:geom-prohibition}
Proposition~\ref{prop:K-audit}(iv) is not a limitation
paragraph---it is a \emph{structure theorem}.  Geometry
identification (claim~C4 in the governance hierarchy of
Section~\ref{sec:Q-status}) is \emph{prohibited}, not merely
difficult, at the scalar-amplitude level.  Breaking this
requires a finer kernel $K'$ incorporating directional
information (Tier~2, Chapter~\ref{ch:future}).
\end{remark}

\begin{remark}[Identifiability gate $G_{\mathrm{id}}$]
\label{rem:G-id}
The current identifiability status is summarised by the triple
$G_{\mathrm{id}} =
(\mathrm{FPR}_{\mathrm{geo}} = 22\%,\;
\mathrm{diag}(\mathbf{C}_{\mathrm{conf}}) = 67\%,\;
\#\text{classes} = 7)$.
Geometry claims require
$\mathrm{FPR}_{\mathrm{geo}} < 10\%$; current status:
\textbf{blocked}.  This gate is an \emph{inference adequacy}
monitor, not a perturbative validity parameter---it tests
source-attribution power, not approximation accuracy.
\end{remark}


% ============================================================
\section{The PSTF multipole expansion and CMB connection}
\label{sec:pstf-cmb}
% ============================================================

The CMB temperature anisotropy, observed along a spatial
direction $e^a$ (with $e_a u^a = 0$, $e_a e^a = 1$), admits
the PSTF multipole expansion
\begin{equation}\label{eq:CMB-multipole}
\frac{\Delta T(\mathbf{e})}{T_0}
= \sum_{\ell=1}^{\infty}
  \eps_\ell\,P_\ell(\cos\vartheta)\,,
\end{equation}
where $\vartheta$ is the angle between $e^a$ and the symmetry
axis $\hat{n}^a$ of the Bianchi model, $P_\ell$ is the Legendre
polynomial, and $\eps_\ell$ is the dimensionless amplitude of
the $\ell$-th multipole.  For axisymmetric Bianchi models, the
$m \neq 0$ modes vanish by symmetry and the anisotropy is fully
characterised by the sequence $\{\eps_\ell\}$.

The relation to the spherical-harmonic coefficients
$a_{\ell m}$ of Chapter~\ref{ch:intro} is
$\eps_\ell = |a_{\ell 0}| \sqrt{4\pi/(2\ell+1)}/T_0$
for $m = 0$.  The first three multipoles---dipole $\eps_1$,
quadrupole $\eps_2$, and octupole $\eps_3$---enter the MES
bound hierarchy derived in
Chapter~\ref{ch:bianchi-bounds} and dominate the constraint
analysis.


% ============================================================
\section{The three-frame problem}\label{sec:frame-problem}
% ============================================================

The passage from a Bianchi geometry to CMB observations involves
three distinct reference frames, each with its own kinematic
decomposition (the conditioned legacy figure gallery in Appendix~\ref{app:conditioned-legacy-figure-gallery}).

The \emph{geometry frame}~(G) is defined by the timelike
eigenvector of the Bianchi symmetry group---the unit normal to
the group orbits.  In this frame the spatial metric is
homogeneous, and the kinematic quantities $\sigma_{ab}$,
$\omega_{ab}$ are properties of the spacetime geometry itself.
Because $u^a$ is hypersurface-orthogonal, we have
$\omega_{ab}^{(G)} = 0$ and, for a comoving fluid,
$\dot{u}_a^{(G)} = 0$.

The \emph{matter frame}~(M) is defined by the four-velocity
$\hat{u}^a$ of the dominant perfect fluid.  When the fluid
is tilted relative to the geometry frame by rapidity $\beta$,
the boost~\eqref{eq:boost-u} connects the two frames.  The
matter-frame kinematic quantities ($\hat{\sigma}_{ab}$,
$\hat{\omega}_{ab}$, $\hat{\dot{u}}_a$) are generically
non-zero even when the geometry-frame quantities vanish.

The \emph{observer frame}~(O) is the CMB rest frame,
obtained from the matter frame by a further peculiar-velocity
boost.  The CMB multipoles $\eps_\ell$ are measured in this
frame.

What does the Bianchi framework constrain?  The defect identity
and the MES bounds are expressed in terms of geometry-frame
quantities.  The CMB multipoles are measured in the observer
frame.  The connection between these two frames---passing
through the matter frame---introduces frame-dependent terms
that must be tracked.  We now derive the relevant transformation
formulas.


\subsection{King--Ellis boost formulae}\label{sec:king-ellis}

King \& Ellis~\cite{KingEllis1973} derived the exact nonlinear
transformation of all kinematic quantities under a Lorentz boost
from the geometry frame to the matter frame.  For a Bianchi
model with $\omega_{ab}^{(G)} = 0$ and
$\dot{u}_a^{(G)} = 0$, the key results are as follows.

The matter-frame expansion is
\begin{equation}\label{eq:KE-Theta}
\hat{\Theta}
= \gamma\,\Theta + \gamma^3\,\sigma_{ab}\,v^a v^b
  + \tilde{D}_a(\gamma\, v^a)\,,
\end{equation}
where $\gamma = \cosh\beta$, $v^a$ is the tilt velocity with
$|v| = \tanh\beta$, and $\tilde{D}_a$ is the spatially projected
covariant derivative.  In the spatially homogeneous setting,
the divergence term $\tilde{D}_a(\gamma v^a)$ involves the
structure constants and does not generically vanish.

The matter-frame \textbf{vorticity} is
\begin{equation}\label{eq:KE-omega}
\hat{\omega}_{ab}
= \gamma^2\,v_{[a}\bigl[\tfrac{1}{3}\Theta\,v_{b]}
  + \sigma_{b]c}\,v^c\bigr]
  + \gamma^2\,\hat{h}_{[a}{}^c\,\tilde{D}_{|c|}
    (\gamma\,v_{b]})\,.
\end{equation}
The leading term $\gamma^2 v_{[a}\sigma_{b]c}v^c$ is the
shear--tilt coupling that generates matter-frame vorticity
from an irrotational geometry frame---a purely relativistic
effect analogous to Thomas precession.  This vorticity vanishes
if and only if either $\sigma_{ab} = 0$ (no shear) or $v^a$
is an eigenvector of $\sigma_{ab}$ (tilt aligned with a
principal shear axis).  For a general tilt direction in a
Bianchi model with non-degenerate shear eigenvalues,
$\hat{\omega}_{ab} \neq 0$.

The matter-frame \textbf{four-acceleration} is
\begin{equation}\label{eq:KE-udot}
\hat{\dot{u}}_a
= \gamma^2\,\hat{h}_a{}^b\bigl[
  \dot{v}_b + \tfrac{1}{3}\Theta\,v_b
  + \sigma_{bc}\,v^c
  + \gamma^2(\sigma_{cd}\,v^c v^d)\,v_b
\bigr]\,,
\end{equation}
where $\dot{v}_b = u^c\nabla_c v_b$ is the geometry-frame time
derivative of the tilt velocity.  King \& Ellis proved the
sharp result: \emph{for a perfect fluid with $p \neq 0$ in an
expanding tilted universe, $\hat{\dot{u}}_a \neq 0$
necessarily}.  Only pressureless dust ($w = 0$) can be
simultaneously tilted and geodesic.

In the \textbf{linearised regime} ($v \ll 1$, $\gamma \approx 1$),
these expressions simplify
to~\cite{EllisVanElst1999,EllisMaartens2012}:
\begin{align}
\hat{\Theta} &\approx \Theta + \tilde{D}^a v_a\,,
\label{eq:KE-lin-Theta}\\
\hat{\sigma}_{ab} &\approx \sigma_{ab}
  + \tilde{D}_{\langle a} v_{b\rangle}\,,
\label{eq:KE-lin-sigma}\\
\hat{\omega}_{ab} &\approx \tilde{D}_{[a}\,v_{b]}\,,
\label{eq:KE-lin-omega}\\
\hat{\dot{u}}_a &\approx \dot{v}_a
  + \tfrac{1}{3}\Theta\,v_a
  + \sigma_{ab}\,v^b\,.
\label{eq:KE-lin-udot}
\end{align}
The linearised acceleration~\eqref{eq:KE-lin-udot} is driven
by three effects: the time evolution of tilt ($\dot{v}_a$),
Hubble drag ($\frac{1}{3}\Theta v_a$), and the shear--tilt
coupling ($\sigma_{ab}v^b$).  For the dipole anomaly amplitude
($\beta \sim 10^{-3}$), the linearised approximation is
accurate to better than $10^{-6}$.


\subsection{The acceleration dipole}\label{sec:accel-dipole}

Our goal in this subsection is to derive the contribution of
the matter-frame acceleration to the CMB dipole---the key
result of VT-07.

The 1+3 covariant Boltzmann equation for the radiation dipole
$I_a^{(1)}$ contains the source
term~\cite{MaartensGE1999,ChallinorLasenby1999}
\begin{equation}\label{eq:boltz-dipole}
\dot{I}_a^{(1)} \ni \tfrac{4}{3}\,I\,\dot{u}_a\,,
\end{equation}
where $I \propto T^4$ is the radiation monopole.  The
acceleration generates a dipolar temperature anisotropy with
amplitude
\begin{equation}\label{eq:eps1-udot-raw}
\eps_1^{(\dot{u})}
\sim \frac{|\hat{\dot{u}}|}{\Theta}\,.
\end{equation}
For a barotropic perfect fluid with equation of state
$\hat{p} = w\hat{\rho}$, the relativistic Euler equation gives
$(\hat{\rho} + \hat{p})\hat{\dot{u}}_a
= -\hat{h}_a{}^b\,\nabla_b \hat{p}$.
In the linearised tilt regime, the dominant contribution
to the acceleration comes from the Hubble drag term in
Eq.~\eqref{eq:KE-lin-udot}.  Specifically, for a radiation-era
fluid ($w = 1/3$), the tilt rapidity $\beta$ is frozen
($\dot{\beta} = 0$; see Section~\ref{sec:tilt-raych}), and
the acceleration reduces to
\begin{equation}
|\hat{\dot{u}}| \approx \Theta\,
  \frac{w}{3(1+w)}\,\beta\,,
\end{equation}
yielding the acceleration efficiency
\begin{equation}\label{eq:eta-udot}
\etaud := \frac{w}{3(1+w)}\,,
\end{equation}
and the acceleration dipole
\begin{equation}\label{eq:eps1-udot}
\eps_1^{(\dot{u})} = \etaud\,\beta\,.
\end{equation}
For radiation ($w = 1/3$), $\etaud = 1/12 \approx 0.083$.
At the matter--radiation equality and in the dust era,
$\etaud \to 0$ (dust is geodesic), so the acceleration
contribution is generated primarily during the radiation epoch
and frozen thereafter.

The total observed dipole is
\begin{equation}\label{eq:eps1-total}
\eps_1 = (1 + \etaud)\,\beta + \eps_1^{\mathrm{int}}\,,
\end{equation}
where $\eps_1^{\mathrm{int}}$ is any genuinely intrinsic
(geometric) dipole.  The $(1 + \etaud)$ factor means that
interpreting the entire $\eps_1$ as tilt (i.e.\ setting
$\beta = \eps_1$) systematically overestimates the tilt
rapidity by a factor $(1 + \etaud)$.


\subsection{Frame-attribution bias}\label{sec:frame-bias}

\begin{definition}[Frame-attribution bias]
\label{def:frame-bias}
The frame-attribution bias is
\begin{equation}\label{eq:frame-bias}
\Bbias := \Bsig^{(\mathrm{obs})} - \Bsig^{(\mathrm{geom})}\,,
\end{equation}
where $\Bsig^{(\mathrm{obs})}$ uses the total observed
$\eps_1$ and $\Bsig^{(\mathrm{geom})}$ uses only the
geometric dipole $\eps_1^{(\mathrm{geom})}
= \eps_1 - \eps_1^{(\dot{u})}$.
\end{definition}

Since $\Bsig = \frac{5}{3}\eps_1 + 3\eps_2 + \frac{3}{7}\eps_3$
is linear in $\eps_1$, the bias is
\begin{equation}
\Bbias = \tfrac{5}{3}\,\eps_1^{(\dot{u})}
= \tfrac{5}{3}\,\etaud\,\beta\,.
\end{equation}
Using $\beta \approx \eps_1/(1+\etaud)$ from the corrected
safe-route (below), the fractional bias is
\begin{equation}
\frac{\Bbias}{\Bsig}
\approx \frac{\etaud}{1 + \etaud}
\approx 7.7\% \quad (w = \tfrac{1}{3})\,.
\end{equation}
This propagates to a $\sim 16\%$ overestimation of $\xmax$ for
shear-dominated types because $\xmax \propto \Bsig^2$.

\subsection{Corrected safe-route tilt bound}
\label{sec:safe-route}
\label{sec:corrected-bound}

Solving Eq.~\eqref{eq:eps1-total} for $\beta$ with
$\eps_1^{\mathrm{int}} = 0$ (the ``safe route''---attributing
the entire dipole to tilt) gives
\begin{equation}\label{eq:beta-sr}
\betasr = \frac{\eps_1}{1 + \etaud}\,,
\end{equation}
which tightens the uncorrected bound $\beta \leq \eps_1$ by
the factor $1/(1 + \etaud) \approx 0.923$ for $w = 1/3$.
At the observed dipole amplitude
$\eps_1 = 1.234 \times 10^{-3}$, the correction shifts the
maximum tilt from $1.234 \times 10^{-3}$ to
$1.140 \times 10^{-3}$---a $7.7\%$ reduction.

This bias is not a correction to be applied and forgotten;
it is a \emph{testable prediction}.  The acceleration content
of the CMB dipole, encoded in $\etaud$, produces a specific
relationship between the intrinsic and observed dipole
amplitudes that can be tested by Euclid~DR1 and SKA Phase~1
measurements of the matter dipole at different
redshifts~\cite{TsagasReview2025}.


% ============================================================
\section{Tilted Raychaudhuri equation in departure variables}
\label{sec:tilt-raych}
% ============================================================

Our goal in this section is to express the dynamical content
of the Raychaudhuri equation in terms of the departure parameter
and its components, revealing how shear, vorticity, tilt, and
curvature interact during the cosmological expansion.  The
derivation follows the analysis of VT-09a.


\subsection{The deceleration parameter}\label{sec:decel-param}

We define the geometry-frame deceleration parameter
\begin{equation}\label{eq:q-def}
q := -1 - \frac{\dot{\Htheta}}{\Htheta^2}\,,
\end{equation}
so that $\dot{\Htheta} = -(1+q)\Htheta^2$.
For an FLRW model with equation of state $w$,
\begin{equation}\label{eq:q-FLRW}
q_{\mathrm{FLRW}} = \tfrac{1}{2}(1+3w)\,\Omega
  - \OL\,.
\end{equation}
Equation~\eqref{eq:q-FLRW} gives $q = 1$ for radiation-dominated FLRW ($w = 1/3$,
$\Omega = 1$) and $q = 1/2$ for matter-dominated FLRW
($w = 0$, $\Omega = 1$).

For a general Bianchi model, the Raychaudhuri
equation~\eqref{eq:raychaudhuri} gives (dividing by
$3\Htheta^2 = \Theta^2/3$ and using the density
parameters)
\begin{equation}\label{eq:raych-normalised}
\frac{\dot{\Theta}}{3\Htheta^2}
= -1 - 2(\Sigstd - \Wstd)
  - \tfrac{1}{2}(1+3w)\,\Omega_{\mathrm{tot}}
  + \OL
  + \frac{\tilde{\nabla}_a\dot{u}^a + |\dot{u}|^2}{3\Htheta^2}\,.
\end{equation}
Rearranging in terms of $q$:
\begin{equation}\label{eq:q-general}
q = q_{\mathrm{FLRW}}
  + 2(\Sigstd - \Wstd)
  + \tfrac{1}{2}(1+3w)\,\Otilt
  - \frac{\tilde{\nabla}_a\dot{u}^a
    + |\dot{u}|^2}{3\Htheta^2}\,,
\end{equation}
where we used the tilt split $\Omega_{\mathrm{tot}}
= \Omega + \Otilt$ and identified the FLRW deceleration
from Eq.~\eqref{eq:q-FLRW}.

Hence the \emph{deceleration excess} due to anisotropy is
\begin{equation}\label{eq:delta-q}
\delta q := q - q_{\mathrm{FLRW}}
= 2\,\Sigstd
  - 2\,\Wstd
  + \tfrac{1}{2}(1+3w)\,\Otilt
  + \delta q_{\mathrm{acc}}\,,
\end{equation}
where $\delta q_{\mathrm{acc}}$ collects the acceleration
terms.  Each term has a transparent physical interpretation:
shear increases deceleration (focusing), vorticity decreases
it (centrifugal support), tilt contributes additional
gravitational deceleration through the enhanced energy
density, and acceleration provides pressure support.


\subsection{Matter-frame deceleration}\label{sec:matter-decel}

For comparison with supernova observations, we need the
deceleration parameter along the matter worldlines.  The
matter-frame Hubble rate $\hat{H}$ differs from the
geometry-frame rate by the tilt correction
$\hat{H} \approx H(1 + \sigma_{ab}v^a v^b/\Theta)$.
The matter-frame deceleration parameter $\qhat$ receives
contributions from both the geometry-frame anisotropy and
the frame transformation itself.

In departure variables, the leading-order result is
\begin{equation}\label{eq:qhat}
\qhat = q_{\mathrm{FLRW}} + \delta q
  + \delta q_{\mathrm{frame}}\,,
\end{equation}
where $\delta q_{\mathrm{frame}}$ contains the
King--Ellis boost corrections to the expansion rate.  For
$\beta \sim 10^{-3}$, the frame correction
$\delta q_{\mathrm{frame}} \sim O(\beta^2) \sim 10^{-6}$
is negligible compared to $\delta q \sim \Sigstd \sim 10^{-6}$.
The matter-frame and geometry-frame deceleration parameters
thus agree to the precision of the present framework.


\subsection{Decay rates and the defect hierarchy}
\label{sec:decay-rates}

The time derivatives of the standardised variables, evaluated
in the near-FLRW regime, reveal the dynamical stability of the
defect decomposition.  From the shear propagation
equation~\eqref{eq:shear-prop}, the leading-order evolution
of $\Sigstd$ is
\begin{equation}\label{eq:Sigstd-dot}
\dot{\Sigstd} \approx -2(1 - q)\,\Htheta\,\Sigstd\,,
\end{equation}
giving $\Sigstd \propto a^{-2(1-q)}$: frozen
($\Sigstd = \mathrm{const}$) in the radiation era ($q = 1$,
since numerator and denominator of
$\Sigstd = \sigma_{ab}\sigma^{ab}/(6\Htheta^2)$ both scale
as~$a^{-4}$), and decaying as $a^{-1}$ in the matter era
($q = 1/2$).

The vorticity decays faster.  From the vorticity propagation
equation~\eqref{eq:vort-prop}, the leading-order result is
\begin{equation}\label{eq:Wstd-dot}
\dot{\Wstd} \approx -2(2 - \tfrac{3w}{2})\,\Htheta\,\Wstd\,,
\end{equation}
giving $\Wstd \propto a^{-2(2-3w/2)}$: decay as $a^{-3}$ for
radiation ($w = 1/3$) and $a^{-4}$ for dust ($w = 0$).  Vorticity
is subdominant to shear at late times for all $w < 2/3$.

The tilt rapidity $\beta$ satisfies the evolution equation
\begin{equation}\label{eq:beta-dot}
\dot{\beta}
= -(1 - 3w)\,\frac{\Htheta}{3}\,\beta
  + (\text{acceleration sources})\,.
\end{equation}
For radiation ($w = 1/3$), the coefficient vanishes:
$\dot{\beta} = 0$.  Tilt is \emph{frozen} during the
radiation era.  For dust ($w = 0$), tilt decays as $\beta
\propto a^{-1/3}$ (ignoring acceleration sources).  The
radiation-era freezing is the physical reason why the tilt
rapidity inferred from the CMB dipole anomaly at $z \approx
1089$ is directly comparable to the present-epoch bulk flow
measured by CF4~\cite{Watkins2023}: tilt survives from the
radiation era without cosmological dilution.

\begin{remark}[Conservation of the defect hierarchy]
\label{rem:hierarchy-conservation}
The decay rates $\Sigstd \gg \Wstd$ at late times
imply that the geodesic shear--vorticity ordering established by the
MES bound coefficients ($\Bsig > \Bom$ under the residual-dipole
convention $\eps_1 < \eps_1^{\mathrm{crit}}$;
Chapter~\ref{ch:bianchi-bounds}) is preserved under dynamical
evolution, not merely an algebraic coincidence.  The
defect-variable formulation makes this stability manifest:
the faster-decaying terms contribute less to $x$ at late times,
and the shear term dominates the defect budget for all
observationally relevant epochs.
\end{remark}


\subsection{Departure evolution equation}\label{sec:defect-evol}

Taking the covariant time derivative of the master
identity~\eqref{eq:master-identity} gives the exact evolution
\begin{equation}\label{eq:xdot}
\dot{x} = \dot{\Sigstd} - \dot{\Wstd}
  + \dot{\Otilt} + \dot{\Okaniso}\,.
\end{equation}
Each term on the right-hand side is governed by the
corresponding propagation equation.  In the near-FLRW regime,
using the decay rates above:
\begin{equation}\label{eq:xdot-approx}
\dot{x} \approx -2(1 - q)\,\Htheta\,\Sigstd
  + 2(2 - \tfrac{3w}{2})\,\Htheta\,\Wstd
  + \dot{\Otilt} + \dot{\Okaniso}\,.
\end{equation}
For the observationally relevant case where $\Sigstd$
dominates $x$, the defect decays at the same rate as the
shear: $x \approx \Sigstd \propto a^{-2(1-q)}$.

Before closing this chapter, we note that the framework
developed here---the kinematic decomposition, the
standardised variables, the master identity, the frame
corrections, the three-layer inference chain, and the defect
evolution---is exact within the Bianchi symmetry class and a
chosen perfect-fluid equation of state.  It is not
``model-independent'' in the sense of holding for arbitrary
matter content: the tilt contribution $\Otilt$ assumes a
perfect fluid with $w > -1$, and the MES ceiling assumes the
Copernican extension postulate.  What the framework does provide
is \emph{assumption-explicit} bookkeeping: every approximation
and every dictionary choice is localised to a specific layer
(algebraic identity, bound application, or observational
connection), and the three-layer inference structure
(Section~\ref{sec:three-layer}) makes the interpretive
dependencies auditable.  The approximations that connect the
exact algebra to observational data enter only in
Chapters~\ref{ch:dipole}--\ref{ch:results}.


% === INLINED: ch02_sec7_tilted_kinematics.tex ===
% Tilted kinematics sections
\section{Tilted Kinematics}
\label{sec:tilted-kinematics}

The kinematics of a tilted cosmology are characterised by the tilt
four-velocity $u^a \neq n^a$, where $n^a$ is the normal to the homogeneous
hypersurfaces.  The tilt rapidity $\beta$ satisfies $\sinh\beta = |u^a n_a|/\sqrt{-n^b n_b}$.

\section{Raychaudhuri Equation and the Defect Decomposition}
\label{sec:tilted-raych-defect}

The Raychaudhuri equation in a tilted Bianchi spacetime acquires
correction terms from the shear and tilt:
\begin{equation}
\dot{\theta} = -\tfrac{1}{3}\theta^2 - 2\sigma^2 + 2\omega^2 + \dot{u}^a{}_{;a} - \tfrac{1}{2}(\rho + 3p) + \Lambda.
\end{equation}
The defect $x = \Sigstd - \Wstd + \Otilt + \Okaniso$ measures the aggregate departure from the FLRW Raychaudhuri relation.

% === END INLINED ===

% === INLINED: ch02_sec8_multi_fluid.tex ===
% ============================================================
% ch02_sec8_multi_fluid.tex
% Multi-fluid species-resolved framework
% Part of: Tetrad-Based Departure Decomposition (Jiwon Park)
% Session S2 draft — 2026-03-30
% ============================================================

\section{Multi-fluid species-resolved framework}
\label{sec:multi-fluid}

Our goal in this section is to extend the single-fluid
departure decomposition developed in
Sections~\ref{sec:std-vars}--\ref{sec:defect-evol} to a
realistic cosmological medium consisting of multiple
interacting species---photons, neutrinos, baryons, cold dark
matter, and dark energy.  The physical motivation is
immediate: the Phase~1.0 Boltzmann solver results of
Chapter~\ref{ch:results} show that neutrinos contribute
$98.4\%$ of the CMB quadrupole power $D_2$
(Section~\ref{sec:sp03-result}), a dominance that is
invisible to a single-fluid treatment.  To connect the
departure parameter $x_C$ to species-resolved observables, we
need a covariant framework that assigns each species its own
kinematic variables while maintaining a single, unambiguous
reference frame for the geometry.

The construction proceeds in five steps.  We first establish
why a single common frame is mandatory
(Section~\ref{sec:common-frame}).  We then classify species
into kinetic and fluid sectors
(Section~\ref{sec:species-class}).  For each kinetic species,
we define a per-species effective-temperature manifold
(Section~\ref{sec:species-teff}).  We derive the macroscopic
reconstruction that translates angular temperature fields
into stress-energy tensors
(Section~\ref{sec:macro-recon}).  We close with the
conservation and interaction bookkeeping that couples
all species through the Einstein equations
(Section~\ref{sec:species-conservation}).


% ============================================================
\subsection{The common frame contract}
\label{sec:common-frame}
% ============================================================

Before diving into the multi-fluid formalism, we must address
a foundational question: in which frame do we define the
kinematic quantities of each species?  Two strategies are
logically possible.  The first assigns each species its own
comoving frame, derives its kinematic quantities there, and
boosts the results back to a common frame at the end.  The
second fixes a single reference congruence $u^a$ from the
outset and describes every species relative to that frame,
encoding each species' bulk motion as a spatial drift vector.
We now show that the second strategy is the only one
consistent with the Bianchi symmetry structure and the
departure decomposition.

\begin{definition}[Common reference congruence]
\label{def:common-frame}
Let $(\mathcal{M}, g_{ab})$ be a Bianchi spacetime with
group-invariant unit normal $n^a$.  The \emph{common reference
congruence} is the choice $u^a = n^a$.  All species-specific
kinematic quantities are defined by projection relative to
this single $u^a$.
\end{definition}

The choice $u^a = n^a$ is not arbitrary.  The standardised
kinematic variables of Section~\ref{sec:std-vars}---the
shear parameter $\Sigstd$, the vorticity parameter $\Wstd$,
and the curvature parameter $\Okaniso$---are defined relative
to the geometry-frame expansion $\Theta$ and shear
$\sigma_{ab}$, both of which belong to the group-invariant
congruence $n^a$.  The master departure identity
(Eq.~\ref{eq:master-identity}) is an algebraic identity in
these geometry-frame variables.  Switching to a
species-comoving frame would replace $\Theta$ and
$\sigma_{ab}$ with tilted counterparts, and the identity
would acquire frame-dependent correction terms that entangle
geometry and matter content in a way that defeats the purpose
of a universal bookkeeping scheme.


\paragraph{Boost pathology for multi-component systems.}

There is a deeper reason to avoid species-comoving frames.
The King--Ellis boost formulae of
Section~\ref{sec:king-ellis} are exact for a single perfect
fluid with known tilt rapidity $\beta$.  For two fluids with
different tilt vectors $v_1^a$ and $v_2^a$, one would need to
boost each fluid separately and then superpose the boosted
stress-energy tensors.  The superposition is algebraically
correct, but the resulting $T^{ab}_{\rm tot}$ is no longer a
single perfect fluid in any frame, so the inverse
problem---recovering species-resolved information from
$T^{ab}_{\rm tot}$---becomes ill-posed.

For kinetic species with nonzero chemical potential
$\mu_s \neq 0$, the situation is worse.  A Lorentz boost
applied to the effective-temperature distribution
$f_s \propto \Phi_{\xi_s}(x_s/\Theta_s - \eta_s)$
mixes the temperature field $\Theta_s(\hat{e})$ and the
fugacity field $\eta_s(\hat{e})$ in a
direction-dependent manner.  That is, if we denote the
boosted quantities by tildes, the transformation
$(\Theta_s, \eta_s) \to (\tilde{\Theta}_s,
\tilde{\eta}_s)$ does not factorise into separate
transformations of $\Theta_s$ and $\eta_s$.  For the
massless $\mu = 0$ case the boost law is clean:
$\tilde{\Theta}_s(\hat{e}') = \Theta_s(\hat{e})
/\gamma(1 - \mathbf{v} \cdot \hat{e})$, which is a pure
rescaling of the temperature field.  Once $\eta_s \neq 0$,
the boosted distribution function can be written
\begin{equation}\label{eq:boost-pathology}
\tilde{f}_s
= \Phi_{\xi_s}\!\left(
  \frac{\tilde{x}_s}{\tilde{\Theta}_s(\hat{e}')}
  - \tilde{\eta}_s(\hat{e}')
\right)
\end{equation}
only if $\tilde{\eta}_s$ absorbs a direction-dependent
shift that depends on both $\beta$ and $\Theta_s$.  The
one-field effective-temperature parameterisation is
therefore not boost-covariant for $\mu \neq 0$ species.

This pathology has a simple resolution: we do not boost.  We
define all species in the common frame $u^a = n^a$ from the
outset, and encode each species' bulk motion through the
spatial components of its angular multipole expansion
(Section~\ref{sec:species-teff}).  The dipole of
$\Theta_s(\hat{e})$ absorbs the species drift, the
quadrupole absorbs the species' response to the background
shear, and so on.  No frame transformation is ever required.


\paragraph{Frame-choice summary.}

We adopt the following contract, valid throughout the
remainder of this work.

\begin{remark}[Common frame contract]
\label{rem:frame-contract}
\emph{(i)}~The geometry-frame normal $n^a$ serves as the
unique reference congruence.
\emph{(ii)}~The spatial projector $h_{ab} = g_{ab} + n_a n_b$
and the kinematic quantities
$(\Theta, \sigma_{ab}, \omega_{ab}, \dot{u}_a)$ are those of
$n^a$, as in Section~\ref{sec:13cov}.
\emph{(iii)}~Each species' bulk velocity relative to $n^a$
is encoded in the dipole $\Theta^{(s)}_a$ of its angular
temperature field, not in a separate Lorentz boost.
\emph{(iv)}~The master departure identity
(Eq.~\ref{eq:master-identity}) and the MES bound hierarchy
(Chapter~\ref{ch:bianchi-bounds}) remain expressed in
geometry-frame variables.
\end{remark}

This contract is consistent with the multi-fluid covariant
formalism of Tsagas, Challinor \&
Maartens~\cite{Tsagas2008}, who decompose
each species' stress-energy relative to a single common
$u^a$ (their~\S2.4).  It is also the natural frame for
the Bianchi Boltzmann solver of Chapter~\ref{ch:results},
which evolves species hierarchies in the group-invariant
tetrad.


% ============================================================
\subsection{Species classification}
\label{sec:species-class}
% ============================================================

Not all cosmological species require the same level of
kinetic description.  A photon gas with its full angular
distribution demands a very different treatment from cold
dark matter, which is adequately described by a pressureless
perfect fluid.  We now establish a classification that
assigns each species to the appropriate tier.

\begin{definition}[Kinetic and fluid sectors]
\label{def:sectors}
A cosmological species $s$ belongs to the \emph{kinetic
sector}~$\mathcal{K}$ if its distribution function $f_s$
carries physically relevant angular structure that must be
tracked.  It belongs to the \emph{fluid
sector}~$\mathcal{F}$ if a finite set of fluid variables
$(\rho_\alpha, p_\alpha, q_\alpha^a, \pi_\alpha^{ab})$
suffices.  The full matter content is
\begin{equation}\label{eq:sector-union}
\{\text{all species}\}
= \mathcal{K} \;\cup\; \mathcal{F}\,,
\qquad
\mathcal{K} \cap \mathcal{F} = \varnothing\,.
\end{equation}
\end{definition}

The classification is not a modelling choice made in
ignorance of the physics.  It follows from the collision
rates and the observational signatures of each species.
We now state the decision criteria and apply them to the
standard $\Lambda$CDM inventory.


\paragraph{Criterion~1: collision rate versus expansion rate.}

A species whose mean free path $\lambda_{\rm mfp}$ satisfies
$\lambda_{\rm mfp} \gg H^{-1}$ at the epoch of interest is
effectively collisionless.  Its distribution function
evolves by free-streaming, and the angular multipole
hierarchy is not truncated by collisions.  Such a species
belongs to $\mathcal{K}$.  Conversely, a species with
$\lambda_{\rm mfp} \ll H^{-1}$ is collision-dominated:
its angular distribution is driven toward local thermodynamic
equilibrium, and a fluid description with low-order
moments suffices.  Such a species belongs to $\mathcal{F}$.


\paragraph{Criterion~2: angular information content.}

Even a collisionless species may be placed in
$\mathcal{F}$ if its angular structure does not contribute
measurably to the observables of interest.  Cold dark
matter, for example, is collisionless after kinetic
decoupling, but its velocity dispersion is negligible
compared to the Hubble flow, and its higher multipoles
($\ell \geq 2$) are suppressed by a factor $(v_{\rm th}/c)^2
\sim 10^{-20}$ relative to the photon quadrupole.  The
physically meaningful test is whether the species'
anisotropic stress $\pi_s^{ab}$ or energy flux $q_s^a$
affects the observables at the target precision.


\paragraph{Criterion~3: number-flux separability.}

For a kinetic species $s \in \mathcal{K}$, a further
subdivision determines whether one or two angular fields are
needed.  If the species carries a conserved quantum number
(lepton number, baryon number) whose directional flux
$j_s^a$ is physically distinct from the energy flux $q_s^a$,
then the effective-temperature ansatz must include both a
temperature field $\Theta_s(\hat{e})$ and a fugacity field
$\eta_s(\hat{e})$.  We call this the \emph{two-field}
configuration.  If the species has vanishing or negligible
chemical potential ($\mu_s / k_B T_s \ll 1$) and its
number flux is locked to its energy flux by
$j_s^a \propto q_s^a$, then a single temperature field
$\Theta_s(\hat{e})$ suffices.  We call this the
\emph{one-field} configuration.

The physical reason for the two-field minimum is a result
from the effective-temperature literature: a scalar
chemical potential $\eta_{0,s}$ does not change the ratio
of number flux to energy flux,
$j_s^a / q_s^a$, so it cannot capture a directional
redistribution of particle number independent of energy.
The direction-dependent fugacity $\eta_s(\hat{e})$ is the
minimal additional degree of freedom that resolves this
degeneracy.

\begin{remark}[One-field diagnostic pitfall]
\label{rem:1d-pitfall}
Applying a one-dimensional tangency diagnostic
$D_{s,\geq 2}$ (Section~\ref{sec:species-tangency}) to a
number-conserving species that genuinely requires two
fields can produce a spuriously large off-manifold
residual.  The excess is not a sign of physical
departure from equilibrium but an artefact of the
under-parameterised manifold: the $\eta$-tangent content
of the collision source is misattributed to higher
energy-shape modes.  For Fermi--Dirac species at moderate
degeneracy, this false-alarm fraction can reach
${\sim}\,98\%$ of the total
$D_{s,\geq 2}$ in the legacy internal stress calculation.
Until a citable dedicated $\Teff$ closure paper or generated
source table is attached, this number is treated as an internal
diagnostic example rather than a public claim.
\end{remark}


\paragraph{Application to $\Lambda$CDM.}

We now apply the three criteria to the standard cosmological
inventory.

\emph{Photons ($\gamma$).}  Before recombination, photons
are tightly coupled to baryons via Thomson scattering
($\lambda_{\rm mfp} \ll H^{-1}$), but the anisotropy
generated by the background shear $\sigma_{ab}$ is the
primary observable.  After recombination, photons
free-stream.  The chemical potential is negligible at the
precision of current CMB measurements ($|\mu_\gamma| /
k_B T_\gamma < 9 \times 10^{-5}$,
Fixsen et al.~\cite{Fixsen1996}).
Hence: $\gamma \in \mathcal{K}$, one-field.

\emph{Neutrinos ($\nu$).}  Neutrinos decouple at
$T \sim 1$~MeV and free-stream thereafter.  Their
anisotropic stress $\pi_\nu^{ab}$ dominates the CMB
quadrupole ($98.4\%$ of $D_2$ at recombination;
Section~\ref{sec:sp03-result}).  The relic neutrino
chemical potential is bounded by Big Bang nucleosynthesis
to $|\eta_\nu| \lesssim 0.07$~\cite{Mangano2012}, which
is small but not negligible for a precision treatment.
Hence: $\nu \in \mathcal{K}$, two-field for maximal
generality (one-field as a controlled approximation when
$\eta_\nu \to 0$).

\emph{Baryons ($b$).}  Before recombination, baryons are
tightly coupled to photons and share their angular
structure.  After recombination, the baryon velocity
dispersion drops to $v_{\rm th} \sim 6$~km/s,
and higher multipoles are negligible.  The baryon--CDM
relative velocity $v_{bc} \sim 30$~km/s
(Tseliakhovich \& Hirata~\cite{TseliakhovichHirata2010})
is coherent on $\sim\!\text{few Mpc}$ scales and affects
small-scale structure formation, but its contribution to
the global departure parameter is
$(v_{bc}/c)^2 \sim 10^{-8}$, far below the sensitivity
of any current survey.  Hence: $b \in \mathcal{F}$.

\emph{Cold dark matter ($c$).}  Collisionless and
pressureless after kinetic decoupling, with
$v_{\rm th}/c \sim 10^{-10}$ at matter--radiation
equality.  No measurable angular structure.
Hence: $c \in \mathcal{F}$.

\emph{Dark energy ($\Lambda$).}  A cosmological constant
has $T^{ab}_\Lambda = -(\Lambda/8\pi G)\,g^{ab}$ with no
kinetic degrees of freedom.  Even for dynamical dark
energy models, the effective fluid description with
$(\rho_{\rm DE}, p_{\rm DE}, w_{\rm DE})$ is standard.
Hence: $\Lambda \in \mathcal{F}$.


\paragraph{Summary.}

For the standard $\Lambda$CDM model, the classification is
\begin{equation}\label{eq:species-partition}
\mathcal{K} = \{\gamma,\; \nu_e,\;
  \bar{\nu}_e,\; \nu_x,\; \bar{\nu}_x\}\,,
\qquad
\mathcal{F} = \{b,\; c,\; \Lambda\}\,,
\end{equation}
where $\nu_x$ denotes the $\mu$- and $\tau$-flavour
neutrinos (treated as a single effective species with
degeneracy $g_{\nu_x} = 4$ in the massless limit).  The
photon sector uses one field; the neutrino sector uses two
fields as the default, reducible to one field when the
chemical potential is set to zero.  We adopt this partition
throughout the remainder of this work.  Extensions
to dark radiation, massive neutrinos, or non-standard
species follow by adding entries to $\mathcal{K}$ or
$\mathcal{F}$ as appropriate.


% ============================================================
\subsection{The effective-temperature manifold for each species}
\label{sec:species-teff}
% ============================================================

For each kinetic species $s \in \mathcal{K}$, we now
construct the finite-dimensional manifold that
parameterises its angular distribution function.  The
construction generalises the single-species $\Teff$ ansatz
of Section~\ref{sec:teff-ansatz} to a multi-species
setting with a common reference frame.

\begin{definition}[Species effective-temperature manifold]
\label{def:species-Teff}
For species $s$ with quantum statistics
$\xi_s \in \{0, -1, +1\}$ (Maxwell--Boltzmann,
Fermi--Dirac, Bose--Einstein respectively), the
\emph{species $\Teff$ manifold} is the set
\begin{equation}\label{eq:Teff-manifold}
\mathcal{M}_s^{\rm Teff}
= \left\{
f_s(x^a, p^b)
= \Phi_{\xi_s}\!\left(
  \frac{x_s}{\Theta_s(\hat{e})}
  - \eta_s(\hat{e})
\right)
\right\}\,,
\end{equation}
where $x_s := E_s / (k_B T_{0,s})$ is the dimensionless
energy, $E_s := -u_a p_s^a > 0$ is the particle energy
measured in the common frame $u^a$,
$\hat{e}^a$ is the spatial propagation direction
($u_a \hat{e}^a = 0$, $\hat{e}_a \hat{e}^a = 1$),
$T_{0,s}$ is the species reference temperature,
$\Theta_s(\hat{e}) > 0$ is the normalised angular
temperature field, $\eta_s(\hat{e})$ is the angular
fugacity field, and $\Phi_\xi(z) := 1/(e^z - \xi)$ is
the equilibrium occupation function.
\end{definition}

For a one-field species, we set $\eta_s(\hat{e}) = 0$
identically, and the manifold reduces to the
single-parameter family
$f_s = \Phi_{\xi_s}(x_s/\Theta_s)$ used in
Section~\ref{sec:teff-ansatz}.

Two aspects of this definition merit comment: the choice
of energy variable, and the angular expansion of
$\Theta_s$ and $\eta_s$.


\paragraph{Energy variable.}

The dimensionless energy $x_s = E_s/(k_B T_{0,s})$ is
defined with respect to the common-frame energy $E_s$,
not the species rest-frame energy.  This choice is correct
for two reasons.  First, the common frame $u^a$ is
fixed by the geometry (Definition~\ref{def:common-frame}),
so $E_s = -u_a p_s^a$ is unambiguous.  Second, the
alternative variable
$y_s = (E_s - \mu_s)/(k_B T_{0,s})$, which absorbs the
chemical potential into the energy coordinate, develops a
direction-dependent domain pathology when
$\mu_s \to \mu_s(\hat{e})$: for certain directions
$\hat{e}$, the variable $y_s$ can become negative at low
energies, and the Laguerre polynomial basis used in the
energy-integrated formalism loses its orthogonality
domain.  The $x_s$-variable avoids this pathology because
$x_s = E_s/(k_B T_{0,s}) > 0$ for all $E_s > 0$,
regardless of the angular structure of $\eta_s$.


\paragraph{Angular expansion.}

The fields $\Theta_s(\hat{e})$ and $\eta_s(\hat{e})$
admit PSTF multipole expansions on the sky sphere of the
common frame:
\begin{equation}\label{eq:Theta-s-PSTF}
\Theta_s(\hat{e})
= 1 + \sum_{\ell \geq 1}
  \Theta^{(s)}_{A_\ell}\,
  \hat{e}^{A_\ell}\,,
\end{equation}
\begin{equation}\label{eq:eta-s-PSTF}
\eta_s(\hat{e})
= \eta_{0,s}
+ \sum_{\ell \geq 1}
  \eta^{(s)}_{A_\ell}\,
  \hat{e}^{A_\ell}\,,
\end{equation}
where $A_\ell = a_1 \cdots a_\ell$ is a PSTF multi-index
and $\hat{e}^{A_\ell}
= \hat{e}^{\langle a_1} \cdots \hat{e}^{a_\ell \rangle}$
is the corresponding PSTF product of direction cosines.
The gauge condition
\begin{equation}\label{eq:Theta-s-gauge}
\langle \Theta_s \rangle_\Omega = 1
\end{equation}
(angular average over the sky sphere) fixes $T_{0,s}$ as
the monopole temperature of species $s$.

The physical content of the low multipoles is transparent.
The dipole $\Theta^{(s)}_a$ encodes the species drift
velocity relative to the common frame: to leading order in
the small-anisotropy expansion,
\begin{equation}\label{eq:dipole-drift}
\Theta^{(s)}_a
\simeq \frac{v^{(s)}_a}{c}\,,
\end{equation}
where $v^{(s)}_a$ is the spatial peculiar velocity of
species $s$ relative to $u^a$.  The quadrupole
$\Theta^{(s)}_{ab}$ encodes the species' response to
the background shear $\sigma_{ab}$, and it is this
quantity that enters the anisotropic stress
$\pi_s^{ab}$ through the moment map
(Section~\ref{sec:macro-recon}).  Higher multipoles
describe the species' angular fine structure, generated
by free-streaming or by mode coupling with the background
geometry.

For the fugacity field, the monopole $\eta_{0,s}$ is
the familiar dimensionless chemical potential
$\mu_s/(k_B T_{0,s})$, and the dipole $\eta^{(s)}_a$
carries the directional number flux that is independent
of the energy flux---precisely the degree of freedom
that distinguishes two-field from one-field species
(Criterion~3 of Section~\ref{sec:species-class}).


\paragraph{Admissibility constraints.}

The distribution function must satisfy $f_s \geq 0$ for
all directions $\hat{e}$ and all energies $x_s > 0$.
For Bose--Einstein species ($\xi_s = +1$), this requires
\begin{equation}\label{eq:BE-admissibility}
\frac{x_s}{\Theta_s(\hat{e})}
- \eta_s(\hat{e}) > 0
\qquad \forall\; x_s > 0,\;
  \hat{e} \in S^2\,,
\end{equation}
which, in the limit $x_s \to 0^+$, implies
$\eta_s(\hat{e}) \leq 0$ for all directions.  For
Fermi--Dirac species ($\xi_s = -1$), there is no sign
constraint on $\eta_s$, but the identification
precision of $\eta_s$ from energy-integrated observables
degrades at large positive degeneracy ($\eta_s \gg 1$)
because the Fermi function saturates and the
entropy-Hessian eigenvalue in the $\eta$-direction
becomes small.  For cosmological neutrinos, the
nucleosynthesis bound $|\eta_\nu| \lesssim 0.07$ places
the system well within the regime where identification
is straightforward.


% ============================================================
\subsection{Macroscopic reconstruction}
\label{sec:macro-recon}
% ============================================================

The effective-temperature manifold of
Definition~\ref{def:species-Teff} parameterises the
microscopic distribution function of each kinetic species.
The Einstein equations, however, require macroscopic
quantities: the species energy density $\rho_s$, pressure
$p_s$, energy flux $q_s^a$, and anisotropic stress
$\pi_s^{ab}$.  We now derive the reconstruction formulae
that extract these quantities from the angular fields
$\Theta_s(\hat{e})$ and $\eta_s(\hat{e})$.


\paragraph{Particle current and stress-energy.}

For species $s \in \mathcal{K}$ with internal degeneracy
$g_s$ (spin states $\times$ particle--antiparticle
factor), the particle four-current and stress-energy
tensor are
\begin{equation}\label{eq:Ns-def}
N_s^a
= \frac{g_s}{(2\pi\hbar)^3}
  \int p_s^a\, f_s\, dP_s\,,
\end{equation}
\begin{equation}\label{eq:Ts-def}
T_s^{ab}
= \frac{g_s}{(2\pi\hbar)^3}
  \int p_s^a\, p_s^b\, f_s\, dP_s\,,
\end{equation}
where $dP_s = d^3 p_s / p_s^0$ is the Lorentz-invariant
momentum-space measure.  For massless species, the
on-shell condition $p_s^a p_{sa} = 0$ gives
$p_s^a = E_s(u^a + \hat{e}^a)$, and the measure
simplifies to $dP_s = E_s\, dE_s\, d\Omega$ in the
common frame, where $d\Omega$ is the solid angle element
on the sky sphere.


\paragraph{1+3 decomposition of the species moments.}

Projecting Eqs.~\eqref{eq:Ns-def}--\eqref{eq:Ts-def}
relative to $u^a$ gives the irreducible decomposition
\begin{equation}\label{eq:Ns-decomp}
N_s^a = n_s\, u^a + j_s^a\,,
\qquad u_a j_s^a = 0\,,
\end{equation}
\begin{equation}\label{eq:Ts-decomp}
T_s^{ab}
= \rho_s\, u^a u^b + p_s\, h^{ab}
  + 2\,u^{(a} q_s^{b)} + \pi_s^{ab}\,,
\end{equation}
with the standard identifications
(cf.\ Eq.~\ref{eq:Tab-general} for the total fluid):
\begin{equation}\label{eq:species-projections}
\begin{aligned}
n_s &= -u_a N_s^a\,,
&\qquad
j_s^{\langle a\rangle} &= h^a{}_b\, N_s^b\,,\\[4pt]
\rho_s &= u_a u_b\, T_s^{ab}\,,
&\qquad
p_s &= \tfrac{1}{3}\, h_{ab}\, T_s^{ab}\,,\\[4pt]
q_s^{\langle a\rangle} &= -h^a{}_c\, u_d\, T_s^{cd}\,,
&\qquad
\pi_s^{\langle ab\rangle}
  &= T_s^{\langle ab \rangle}\,.
\end{aligned}
\end{equation}
The number density $n_s$ and number drift $j_s^a$
belong to the particle-current sector; the energy
density $\rho_s$, pressure $p_s$, energy flux $q_s^a$,
and anisotropic stress $\pi_s^{ab}$ belong to the
stress-energy sector.  Both sectors are built from the
same underlying distribution function $f_s$ and hence
from the same angular fields $(\Theta_s, \eta_s)$.


\paragraph{Explicit reconstruction for massless species.}

For a massless species with $p_s^a = E_s(u^a +
\hat{e}^a)$, we substitute the $\Teff$ ansatz
(Eq.~\ref{eq:Teff-manifold}) into the moment integrals.
The energy integration produces standard thermodynamic
integrals, and the angular integration projects out the
PSTF multipoles of $\Theta_s$ and $\eta_s$.

The energy density is
\begin{equation}\label{eq:rho-s-recon}
\rho_s
= \frac{g_s}{2\pi^2\hbar^3}\,(k_B T_{0,s})^4
  \left\langle \Theta_s^4\, \mathcal{I}_3(\xi_s,
  \eta_s) \right\rangle_\Omega\,,
\end{equation}
where $\mathcal{I}_n(\xi, \eta)$ is the generalised
thermodynamic integral
\begin{equation}\label{eq:thermo-integral}
\mathcal{I}_n(\xi, \eta)
:= \int_0^\infty
  \frac{x^n\, dx}{e^{x - \eta} - \xi}\,,
\end{equation}
and the angular brackets $\langle \cdot
\rangle_\Omega$ denote the average over the sky sphere.
For $\xi = +1$ and $\eta = 0$ (massless bosons),
$\mathcal{I}_3(+1, 0) = \pi^4/15$, recovering the
Stefan--Boltzmann law.  For $\xi = -1$ and $\eta = 0$
(massless fermions),
$\mathcal{I}_3(-1, 0) = 7\pi^4/120$.

The pressure follows from the trace of the spatial
stress:
\begin{equation}\label{eq:p-s-recon}
p_s = \frac{1}{3}\,\rho_s\,,
\end{equation}
which is the equation of state for massless
radiation, valid direction by direction before angular
averaging.  This relation is an exact consequence of the
on-shell condition $p_s^a p_{sa} = 0$ and does not
depend on the $\Teff$ ansatz.

The energy flux is
\begin{equation}\label{eq:q-s-recon}
q_s^a
= \frac{g_s}{2\pi^2\hbar^3}\,(k_B T_{0,s})^4
  \left\langle \hat{e}^a\,\Theta_s^4\,
  \mathcal{I}_3(\xi_s, \eta_s)
  \right\rangle_\Omega\,.
\end{equation}
To leading order in the multipole expansion, the
angular average picks out the dipole
$\Theta^{(s)}_a$:
\begin{equation}\label{eq:q-leading}
q_s^a
\simeq \frac{4}{3}\,\rho_s^{(0)}\,
  \Theta^{(s)}_a\,,
\end{equation}
where $\rho_s^{(0)}$ is the isotropic (monopole)
energy density.  This confirms the physical
interpretation of $\Theta^{(s)}_a$ as a normalised
drift velocity (Eq.~\ref{eq:dipole-drift}).

The anisotropic stress is
\begin{equation}\label{eq:pi-s-recon}
\pi_s^{ab}
= \frac{g_s}{2\pi^2\hbar^3}\,(k_B T_{0,s})^4
  \left\langle \hat{e}^{\langle a}\hat{e}^{b\rangle}\,
  \Theta_s^4\, \mathcal{I}_3(\xi_s, \eta_s)
  \right\rangle_\Omega\,,
\end{equation}
and to leading order,
\begin{equation}\label{eq:pi-leading}
\pi_s^{ab}
\simeq \frac{8}{15}\,\rho_s^{(0)}\,
  \Theta^{(s)}_{ab}\,.
\end{equation}
The quadrupole $\Theta^{(s)}_{ab}$ therefore directly
sources the anisotropic stress that enters the
Einstein equations, the shear propagation equation
(Section~\ref{sec:shear-prop}), and the gravitational
Silk damping of the CMB.


\paragraph{Number sector.}

The number density and number drift are reconstructed
from lower-weight integrals:
\begin{equation}\label{eq:n-s-recon}
n_s
= \frac{g_s}{2\pi^2\hbar^3}\,(k_B T_{0,s})^3
  \left\langle \Theta_s^3\,
  \mathcal{I}_2(\xi_s, \eta_s)
  \right\rangle_\Omega\,,
\end{equation}
\begin{equation}\label{eq:j-s-recon}
j_s^a
= \frac{g_s}{2\pi^2\hbar^3}\,(k_B T_{0,s})^3
  \left\langle \hat{e}^a\, \Theta_s^3\,
  \mathcal{I}_2(\xi_s, \eta_s)
  \right\rangle_\Omega\,.
\end{equation}
These integrals involve $\mathcal{I}_2$ rather than
$\mathcal{I}_3$, reflecting the different energy
weight of the number current ($\sim E^2$) compared to
the stress-energy ($\sim E^3$).  For a two-field
species, the dependence of $\mathcal{I}_2(\xi_s,
\eta_s(\hat{e}))$ on the angular fugacity field
$\eta_s(\hat{e})$ is what allows the number drift
$j_s^a$ to vary independently of the energy flux
$q_s^a$.  This independence is the physical content of
Criterion~3 and the mathematical reason why two-field
species require two angular fields.

\begin{proposition}[Number--energy flux independence]
\label{prop:flux-independence}
For a two-field species $s$ with $\eta_s(\hat{e})
\neq \text{const}$, the number drift $j_s^a$ and the
energy flux $q_s^a$ are generically linearly
independent: there exists no scalar function
$\lambda$ such that $j_s^a = \lambda\, q_s^a$ for all
configurations $(\Theta_s, \eta_s)$.
\end{proposition}

\begin{proof}
The energy flux $q_s^a$ is proportional to the angular
average $\langle \hat{e}^a\, \Theta_s^4\,
\mathcal{I}_3 \rangle_\Omega$
(Eq.~\ref{eq:q-s-recon}), while the number drift
$j_s^a$ is proportional to
$\langle \hat{e}^a\, \Theta_s^3\,
\mathcal{I}_2 \rangle_\Omega$
(Eq.~\ref{eq:j-s-recon}).  Consider a configuration
with $\Theta_s = 1 + \Theta^{(s)}_b \hat{e}^b$ and a
dipolar fugacity $\eta_s(\hat{e}) = \eta_{0,s} +
\eta^{(s)}_a \hat{e}^a$.  At linear order, both fluxes
receive contributions from both $\Theta^{(s)}_a$ and
$\eta^{(s)}_a$:
$q_s^a \propto A_E\,\Theta^{(s)}_a
+ B_E\,\eta^{(s)}_a$ and
$j_s^a \propto A_N\,\Theta^{(s)}_a
+ B_N\,\eta^{(s)}_a$,
where $A_E \propto \mathcal{I}_3$,
$B_E \propto \partial\mathcal{I}_3/\partial\eta$,
$A_N \propto \mathcal{I}_2$, and
$B_N \propto \partial\mathcal{I}_2/\partial\eta$.
The proportionality $j_s^a = \lambda\,q_s^a$ would
require $A_N/A_E = B_N/B_E$, i.e.\
$\mathcal{I}_2/\mathcal{I}_3
= (\partial\mathcal{I}_2/\partial\eta)
 /(\partial\mathcal{I}_3/\partial\eta)$.
This ratio identity fails for generic
$(\xi_s, \eta_{0,s})$: differentiation with respect to
$\eta$ shifts the energy weight from $x^n$ to
$x^n e^{x-\eta}/(e^{x-\eta}-\xi)^2$, and the
resulting $\eta$-dependence of
$\partial\mathcal{I}_n/\partial\eta$ differs between
$n = 2$ and $n = 3$.  Hence the $2 \times 2$ matrix
$(A_E, B_E; A_N, B_N)$ generically has rank~2, and
$j_s^a \neq \lambda\,q_s^a$.
\end{proof}


\paragraph{Fluid-sector species.}

For species $\alpha \in \mathcal{F}$, no
distribution-function reconstruction is needed.  The
stress-energy tensor is specified directly by the fluid
variables:
\begin{equation}\label{eq:fluid-Tab}
T_\alpha^{ab}
= \rho_\alpha\, u^a u^b
  + p_\alpha\, h^{ab}
  + 2\,u^{(a} q_\alpha^{b)}
  + \pi_\alpha^{ab}\,,
\end{equation}
with the equation of state $p_\alpha = w_\alpha
\rho_\alpha$ (where $w_b = 0$ for baryons,
$w_c = 0$ for CDM, $w_\Lambda = -1$ for a cosmological
constant) and with $q_\alpha^a$ and $\pi_\alpha^{ab}$
specified by the relevant transport equation---typically
the Euler equation for baryons and CDM, or set to zero
for $\Lambda$.


% ============================================================
\subsection{Conservation and interaction bookkeeping}
\label{sec:species-conservation}
% ============================================================

The species-resolved stress-energy tensors are not
separately conserved: energy and momentum are exchanged
between species through collisions.  We now write the
conservation equations that track these exchanges and
impose the constraint that total energy--momentum is
conserved.


\paragraph{Species transport equations.}

Each kinetic species $s \in \mathcal{K}$ satisfies the
covariant Boltzmann equation
\begin{equation}\label{eq:Boltzmann-s}
\mathcal{L}_s[f_s] = C_s[\{f_r\}]\,,
\end{equation}
where $\mathcal{L}_s$ is the Liouville operator on the
mass shell of species $s$
(cf.\ Section~\ref{sec:liouville}), and $C_s[\{f_r\}]$
is the collision integral, which depends on the
distribution functions of all species $r$ that interact
with $s$.  Integrating Eq.~\eqref{eq:Boltzmann-s} over
momentum space yields the macroscopic conservation
equations:
\begin{equation}\label{eq:number-conservation}
\nabla_a N_s^a = R_s\,,
\end{equation}
\begin{equation}\label{eq:energy-conservation}
\nabla_a T_s^{ab} = Q_s^{\;b}\,,
\end{equation}
where $R_s$ is the net particle creation rate
(vanishing for number-conserving species) and
$Q_s^{\;b}$ is the energy--momentum exchange four-vector
between species $s$ and the rest of the medium.

For each fluid species $\alpha \in \mathcal{F}$, the
analogous equations are
\begin{equation}\label{eq:fluid-conservation}
\nabla_a T_\alpha^{ab} = Q_\alpha^{\;b}\,.
\end{equation}


\paragraph{Global conservation constraint.}

Total energy--momentum conservation,
$\nabla_a T^{ab}_{\rm tot} = 0$, requires
\begin{equation}\label{eq:total-conservation}
\sum_{s \in \mathcal{K}} Q_s^{\;b}
+ \sum_{\alpha \in \mathcal{F}} Q_\alpha^{\;b}
= 0\,.
\end{equation}
This constraint is not an additional equation of motion but an
algebraic constraint on the interaction sources $Q_s^b$
and $Q_\alpha^b$.  For the standard $\Lambda$CDM
inventory, the dominant interaction is Thomson scattering
between photons and baryons,
\begin{equation}\label{eq:Thomson-Q}
Q_\gamma^{\;a} = -Q_b^{\;a}
= -n_e \sigma_T\, \big[
  q_\gamma^a - \tfrac{4}{3}\rho_\gamma\, v_b^a
  + \text{(polarisation terms)}
\big]\,,
\end{equation}
where $n_e$ is the free-electron number density,
$\sigma_T$ is the Thomson cross-section, and $v_b^a$ is
the baryon velocity.  Neutrinos decouple before
recombination and have $Q_\nu^a = 0$ at all epochs of
interest.  Dark matter interacts only gravitationally:
$Q_c^a = 0$.

For a subset $\mathcal{C} \subseteq \mathcal{K}$ of
number-conserving species (e.g.\ electron-type neutrinos
and antineutrinos jointly conserving lepton number), the
number sources satisfy
\begin{equation}\label{eq:number-constraint}
\sum_{s \in \mathcal{C}} R_s = 0\,.
\end{equation}
This bookkeeping is the natural habitat of the two-field
fugacity $\eta_s(\hat{e})$: the directional number
redistribution that $\eta_s$ encodes is driven by
precisely the interactions that conserve total number
within $\mathcal{C}$ while transferring particles between
directions.


\paragraph{Einstein coupling.}

The total stress-energy entering the Einstein equations is
\begin{equation}\label{eq:total-Tab}
T^{ab}_{\rm tot}
= \sum_{s \in \mathcal{K}} T_s^{ab}
+ \sum_{\alpha \in \mathcal{F}} T_\alpha^{ab}\,,
\end{equation}
and the Einstein equation reads
\begin{equation}\label{eq:Einstein-multi}
G^{ab} = 8\pi G\, T^{ab}_{\rm tot}\,,
\end{equation}
where we work in units with $c = 1$.  The geometry
does not see the individual species directly; it couples
to them only through the total reconstructed
stress-energy.  The species-resolved information is
carried by the angular fields $(\Theta_s, \eta_s)$ and
the fluid variables $(\rho_\alpha, p_\alpha,
q_\alpha^a, \pi_\alpha^{ab})$, which enter
$T^{ab}_{\rm tot}$ through the reconstruction formulae
of Section~\ref{sec:macro-recon}.


\paragraph{FLRW limit.}

As a consistency check, we verify the isotropic limit.
Setting $\Theta_s(\hat{e}) = 1$ and
$\eta_s(\hat{e}) = \eta_{0,s}$ (constant) for all
$s \in \mathcal{K}$, and $q_\alpha^a = \pi_\alpha^{ab}
= 0$ for all $\alpha \in \mathcal{F}$, the species
angular structure vanishes:
$j_s^a = q_s^a = \pi_s^{ab} = 0$ for each species.
The total stress-energy reduces to
\begin{equation}\label{eq:FLRW-limit-multi}
T^{ab}_{\rm tot}
= \rho_{\rm tot}\, u^a u^b
  + p_{\rm tot}\, h^{ab}\,,
\end{equation}
with $\rho_{\rm tot} = \sum_s \rho_s + \sum_\alpha
\rho_\alpha$ and $p_{\rm tot} = \sum_s p_s + \sum_\alpha
p_\alpha$.  The Friedmann equation
(Section~\ref{sec:gen-friedmann}) is recovered exactly,
and the departure parameter $x_C = 0$ by the master
identity (Eq.~\ref{eq:master-identity}).  This recovery is the
multi-species generalisation of the
single-fluid FLRW limit derived in
Section~\ref{sec:defect-x}.


\paragraph{Scope of the present formalism.}

Before closing this section, we note explicitly what the
multi-fluid framework does and does not provide.  It
provides a species-resolved, covariant parameterisation
of the matter content (the $\Teff$ manifold per species),
a reconstruction map from angular fields to stress-energy
(the moment integrals), and a bookkeeping structure for
inter-species energy--momentum exchange (the $Q_s^b$
sources).  It does \emph{not} provide a closed dynamical
system: the evolution equations for $\Theta_s(\hat{e})$
and $\eta_s(\hat{e})$ require projection of the
Boltzmann equation onto the $\Teff$ manifold, which is
the subject of Chapter~\ref{ch:teff}
(Sections~\ref{sec:teff-ode}--\ref{sec:defect-prop}).
It also does not provide a global-in-time existence or
stability theorem for the coupled multi-species system;
this is an open mathematical problem even for the
two-fluid super-radiative case
(Fournodavlos, Marshall \&
Oliynyk~\cite{FournodavlosMarshallOliynyk2025};
see Section~\ref{sec:open-problems}).
The framework is best understood as a \emph{monitored
closure layer}: it provides a finite-dimensional
reduction of the infinite-dimensional kinetic
description, together with a per-species diagnostic
(Section~\ref{sec:species-tangency}) that certifies when
the reduction is adequate and flags when it must be
upgraded.


% ============================================================
\subsection{Species-resolved tangency diagnostics}
\label{sec:species-tangency}
% ============================================================

The $\Teff$ manifold of
Definition~\ref{def:species-Teff} parameterises each
kinetic species' distribution function by one or two
angular fields.  We now construct a diagnostic that
measures, for each species independently, whether the
exact collision source is tangent to the species' $\Teff$
manifold.  When the diagnostic is small, the
finite-dimensional parameterisation is adequate; when it
is large, the parameterisation must be upgraded by adding
energy-shape modes.  The construction generalises the
single-species tangency diagnostic of
Section~\ref{sec:tangency} to the multi-species setting.


\paragraph{Per-ray tangent space.}

Fix a spatial direction $\hat{e}$ and a kinetic species
$s \in \mathcal{K}$ with two-field ansatz
$f_s = \Phi_{\xi_s}(x_s/\Theta_s - \eta_s)$.  The
tangent directions to the $\Teff$ manifold at the current
state are obtained by differentiating $f_s$ with respect
to the two parameters $\Theta_s$ and $\eta_s$:
\begin{equation}\label{eq:tangent-dTheta}
\frac{\partial f_s}{\partial \Theta_s}
= w_{*,s}(x_s)\,
  \frac{x_s}{\Theta_s^2}\,,
\end{equation}
\begin{equation}\label{eq:tangent-deta}
\frac{\partial f_s}{\partial \eta_s}
= w_{*,s}(x_s)\,,
\end{equation}
where the statistical weight function is
\begin{equation}\label{eq:stat-weight}
w_{*,s}(x_s)
:= f_s(1 + \xi_s f_s)\,.
\end{equation}
The two tangent directions are proportional to the
functions $x_s$ and $1$ (constant), respectively,
weighted by $w_{*,s}$.  The per-ray tangent space is
therefore
\begin{equation}\label{eq:tangent-space}
V_s = \mathrm{span}\{x_s,\; 1\}\,.
\end{equation}
For a one-field species ($\eta_s = 0$ fixed), the
tangent space reduces to $V_s = \mathrm{span}\{x_s\}$.


\paragraph{The Fisher--entropy inner product.}

The natural inner product on functions of $x_s$ is the
weighted $L^2$ product
\begin{equation}\label{eq:fisher-inner}
\langle g, h \rangle_{*,s}
:= \int_0^\infty
  g(x_s)\, h(x_s)\,
  w_{*,s}(x_s)\, x_s^3\, dx_s\,,
\end{equation}
where the factor $x_s^3$ arises from the
momentum-space measure for massless particles.  This
inner product plays a triple role: it is the
Fisher information metric on the $\Teff$ manifold, the
Hessian of the entropy functional restricted to the
manifold, and the kernel of the orthogonal projection
onto $V_s$.  The triple identification is a structural
specified structural result of the effective-temperature
formalism; its publication-facing status remains conditional on
the generated closure notes and current artifact manifests.


\paragraph{Gram matrix.}

The Gram matrix of $V_s$ with respect to the inner
product~\eqref{eq:fisher-inner} is the $2 \times 2$
symmetric positive-definite matrix
\begin{equation}\label{eq:Gram-matrix}
(\mathcal{G}_s)_{ij}
= \langle v_i,\, v_j \rangle_{*,s}\,,
\qquad
v_1 = x_s,\;\; v_2 = 1\,,
\end{equation}
with entries
\begin{equation}\label{eq:Gram-entries}
\mathcal{G}_s
= \begin{pmatrix}
\langle x_s^2 \rangle_{*,s}
  & \langle x_s \rangle_{*,s} \\[4pt]
\langle x_s \rangle_{*,s}
  & \langle 1 \rangle_{*,s}
\end{pmatrix}\,,
\end{equation}
where we write $\langle x_s^k \rangle_{*,s}
:= \int_0^\infty x_s^{k+3}\, w_{*,s}\, dx_s$ for
brevity.  The positive-definiteness of $\mathcal{G}_s$
is equivalent to the identifiability of $(\Theta_s,
\eta_s)$ from the energy and number sectors: if the
Gram matrix were degenerate, the two parameters could
not be separately reconstructed from observables.  For
Fermi--Dirac species at large positive degeneracy
($\eta_s \gg 1$), the smaller eigenvalue of
$\mathcal{G}_s$ decreases, reflecting the reduced
sensitivity of the Fermi function to $\eta_s$ when the
occupation is near saturation.  For cosmological
neutrinos with $|\eta_\nu| \lesssim 0.07$, both
eigenvalues are of order unity, and identification is
well-conditioned.


\paragraph{Collision field and tangency criterion.}

The collision integral $C_s[\{f_r\}]$
(Eq.~\ref{eq:Boltzmann-s}) drives the distribution
function away from or toward the $\Teff$ manifold.  To
quantify this, we define the \emph{normalised collision
field}
\begin{equation}\label{eq:collision-field}
G_{\xi,s}(x_s)
:= \frac{C_s[\{f_r\}](x_s)}{w_{*,s}(x_s)}\,.
\end{equation}
The division by $w_{*,s}$ converts the collision
integral from a rate of change of $f_s$ into a rate of
change of the ``natural coordinates'' on the manifold.
The tangency condition is:
\begin{equation}\label{eq:tangency-criterion}
G_{\xi,s} \in V_s
\quad \Longleftrightarrow \quad
\text{the collision source for species $s$ is tangent
to its $\Teff$ manifold.}
\end{equation}
When this condition holds, the collision integral
drives $f_s$ along the manifold---changing $\Theta_s$
and $\eta_s$ without exciting higher energy-shape
modes.  When it fails, the collision integral has a
component orthogonal to $V_s$, and the
finite-dimensional parameterisation incurs an error
that grows with time.


\paragraph{The departure diagnostic $D_{s,\geq 2}$.}

The off-manifold content of the collision source is
measured by the orthogonal projection residual:
\begin{equation}\label{eq:D-diagnostic}
D_{s,\geq 2}^2
:= \big\|
  G_{\xi,s} - \Pi_{V_s} G_{\xi,s}
\big\|_{*,s}^2\,,
\end{equation}
where $\Pi_{V_s}$ is the orthogonal projector onto
$V_s$ with respect to the inner
product~\eqref{eq:fisher-inner}.  The subscript
``$\geq 2$'' indicates that the residual lives in the
energy-shape modes of order two and higher---the
components of $G_{\xi,s}$ that cannot be absorbed by
adjusting $\Theta_s$ and $\eta_s$.

For a one-field species, the projection is onto
$V_s = \mathrm{span}\{x_s\}$, and the diagnostic
reduces to
\begin{equation}\label{eq:D-one-field}
D_{s,\geq 2}^{\rm (1D)}
= \left\|
  G_{\xi,s}
  - \frac{\langle x_s,\, G_{\xi,s}\rangle_{*,s}}
         {\langle x_s^2 \rangle_{*,s}}\, x_s
\right\|_{*,s}\,.
\end{equation}


\paragraph{Upgrade policy.}

The diagnostic $D_{s,\geq 2}$ is evaluated per
species and per ray.  It defines a graded response:
\begin{equation}\label{eq:upgrade-policy}
\begin{cases}
D_{s,\geq 2} \ll 1:
  & \text{current $\Teff$ manifold adequate;
    retain $(\Theta_s, \eta_s)$}\,, \\[4pt]
D_{s,\geq 2} \gtrsim \varepsilon_{\rm thr}:
  & \text{add energy-shape mode(s) to species $s$,
    or switch to multigroup.}
\end{cases}
\end{equation}
The threshold $\varepsilon_{\rm thr}$ is set by the
target observable accuracy; for the scalar-amplitude
pipeline of this work, the dominant error sources are
statistical (Section~\ref{sec:combined}), and
$\varepsilon_{\rm thr} \sim 10^{-3}$ is conservative.
The ``diagnose-to-upgrade'' philosophy---monitor
tangency, then add complexity only where and when
needed---is the operational core of the $\Teff$
approach.  It replaces the traditional strategy of
choosing $\ell_{\rm max}$ a priori and hoping that the
truncation error is small.


% ============================================================
\subsection{Projected evolution and the Bianchi hierarchy}
\label{sec:projected-evolution}
% ============================================================

The tangency diagnostic of
Section~\ref{sec:species-tangency} certifies the
instantaneous quality of the $\Teff$ reduction.  To
evolve the species angular fields in time, we need a
dynamical system.  We now derive the projected evolution
equations by combining the covariant Boltzmann equation
(Eq.~\ref{eq:Boltzmann-s}) with the $\Teff$ ansatz and
the macroscopic reconstruction of
Section~\ref{sec:macro-recon}.


\paragraph{State vector and observable vector.}

For species $s$ with a $(L_\Theta, L_\eta)$ truncation
(angular multipoles retained up to $\ell = L_\Theta$
for the temperature field and $\ell = L_\eta$ for the
fugacity field), the state vector is
\begin{equation}\label{eq:state-vector}
\alpha_s
= \Big(
\ln T_{0,s},\;
\eta_{0,s},\;
\Theta^{(s)}_{\langle A_1 \rangle},\;
\ldots,\;
\Theta^{(s)}_{\langle A_{L_\Theta} \rangle},\;
\eta^{(s)}_{\langle A_1 \rangle},\;
\ldots,\;
\eta^{(s)}_{\langle A_{L_\eta} \rangle}
\Big)\,.
\end{equation}
The corresponding observable vector pairs energy-weighted
and number-weighted angular moments:
\begin{equation}\label{eq:obs-vector}
m_s
= \Big(
\mathcal{E}_s,\;
\mathcal{E}^{(s)}_{\langle A_1 \rangle},\;
\ldots,\;
\mathcal{E}^{(s)}_{\langle A_{L_\Theta} \rangle},\;
\mathcal{N}_s,\;
\mathcal{N}^{(s)}_{\langle A_1 \rangle},\;
\ldots,\;
\mathcal{N}^{(s)}_{\langle A_{L_\eta} \rangle}
\Big)\,,
\end{equation}
where
\begin{equation}\label{eq:energy-moments}
\mathcal{E}^{(s)}_{\langle A_\ell \rangle}
:= \Delta_\ell\,
  \frac{g_s\,(k_B T_{0,s})^4}{(2\pi\hbar)^3}
  \int_0^\infty dx_s\, x_s^3\,
  F^{(s)}_{A_\ell}(x_s)
\end{equation}
are the energy-weighted multipoles and
\begin{equation}\label{eq:number-moments}
\mathcal{N}^{(s)}_{\langle A_\ell \rangle}
:= \Delta_\ell\,
  \frac{g_s\,(k_B T_{0,s})^3}{(2\pi\hbar)^3}
  \int_0^\infty dx_s\, x_s^2\,
  F^{(s)}_{A_\ell}(x_s)
\end{equation}
are the number-weighted multipoles.  Here
$F^{(s)}_{A_\ell}(x_s)$ is the $\ell$-th PSTF
coefficient in the angular expansion of $f_s$, and
$\Delta_\ell$ is a normalisation constant.  The $x_s^3$
versus $x_s^2$ weighting distinguishes the energy and
number sectors---the same structural difference that
underlies
Proposition~\ref{prop:flux-independence} and the
two-field requirement for number-conserving species.


\paragraph{The projected evolution system.}

The observable vector $m_s$ depends on the state vector
$\alpha_s$ through the $\Teff$ ansatz and the moment
integrals.  By the chain rule,
\begin{equation}\label{eq:chain-rule}
\dot{m}_s
= J_s(\alpha_s)\, \dot{\alpha}_s\,,
\qquad
J_s := \frac{\partial m_s}{\partial \alpha_s}\,,
\end{equation}
where $J_s$ is the state-to-observable Jacobian.
Substituting the Boltzmann
equation~\eqref{eq:Boltzmann-s} and projecting its
energy-weighted and number-weighted moments gives the
exact reduced system
\begin{equation}\label{eq:projected-system}
\boxed{
J_s(\alpha_s)\, \dot{\alpha}_s
= \mathcal{R}_s[\alpha_s;\, u^a, \Theta,
  \sigma_{ab}, \omega_{ab}, \dot{u}_a, D_a]
+ \mathcal{C}_s[\alpha_s, \{\alpha_r\}_{r \neq s}]\,,
}
\end{equation}
where $\mathcal{R}_s$ is the streaming/kinematic
projection (containing the expansion, shear, vorticity,
acceleration, and spatial-gradient couplings) and
$\mathcal{C}_s$ is the collision projection
(containing the inter-species energy--momentum
exchange).  The system~\eqref{eq:projected-system} is
the species-resolved, two-field generalisation of the
mass-matrix evolution equation derived for the
single-species case in Section~\ref{sec:ode-system}.

The Jacobian $J_s$ is generically a dense, nonlinear
matrix: its entries are thermodynamic integrals that
depend on the current state $(\Theta_s, \eta_s)$.
Near the isotropic state, $J_s$ block-diagonalises by
tensor rank (scalar, vector, tensor blocks decouple at
linear order), and each block reduces to a small
matrix whose entries are standard thermodynamic
functions of $\eta_{0,s}$.  The invertibility of $J_s$
is guaranteed by the positive-definiteness of the
Gram matrix (Eq.~\ref{eq:Gram-matrix}), which ensures
that the state vector can be recovered from the
observable vector.


\paragraph{FLRW isotropic limit.}

Setting all anisotropic multipoles to zero
($\Theta^{(s)}_{A_\ell} = \eta^{(s)}_{A_\ell} = 0$
for $\ell \geq 1$) and all Bianchi kinematic
quantities to their FLRW values
($\sigma_{ab} = \omega_{ab} = \dot{u}_a = D_a = 0$,
$\Theta = 3H$), the system~\eqref{eq:projected-system}
collapses to the $2 \times 2$ monopole system
\begin{equation}\label{eq:FLRW-monopole}
\begin{pmatrix}
3\,n_s & \chi^{(s)}_{N\eta} \\[4pt]
4\,\rho_s & \chi^{(s)}_{E\eta}
\end{pmatrix}
\begin{pmatrix}
\dot{(\ln T_{0,s})} \\[4pt]
\dot{\eta}_{0,s}
\end{pmatrix}
=
\begin{pmatrix}
R_s - 3H\, n_s \\[4pt]
S_{E,s} - 4H\, \rho_s
\end{pmatrix}\,,
\end{equation}
where the susceptibilities are
$\chi^{(s)}_{N\eta}
:= \partial n_s / \partial \eta_{0,s}$
and
$\chi^{(s)}_{E\eta}
:= \partial \rho_s / \partial \eta_{0,s}$,
the number source is $R_s$
(Eq.~\ref{eq:number-conservation}), and the energy
source is $S_{E,s} := -u_b Q_s^{\;b}$
(Eq.~\ref{eq:energy-conservation}).  The factors
$3\,n_s$ and $4\,\rho_s$ on the left-hand side arise
from $\partial_{\ln T_0} n \propto T_0^3$ and
$\partial_{\ln T_0} \rho \propto T_0^4$ for massless
species.

For a collisionless massless species ($R_s = 0$,
$S_{E,s} = 0$), the
system~\eqref{eq:FLRW-monopole} yields
\begin{equation}\label{eq:FLRW-solution}
T_{0,s} \propto a^{-1}\,,
\qquad
\eta_{0,s} = \text{const}\,,
\end{equation}
recovering the standard adiabatic cooling law and the
conservation of the dimensionless chemical potential for
free-streaming relativistic species.  For photons with
rapid number-changing processes ($R_\gamma \neq 0$), the
system drives $\eta_{0,\gamma} \to 0$, consistent with
the observed blackbody spectrum.  For decoupled neutrinos,
$\eta_{0,\nu}$ freezes at its pre-decoupling value, in
agreement with Big Bang nucleosynthesis constraints.


\paragraph{Homogeneous Bianchi limit.}

In a spatially homogeneous Bianchi spacetime, the
spatial gradient $D_a$ vanishes, eliminating the
$\ell \leftrightarrow \ell \pm 1$ coupling that
connects adjacent multipoles in the inhomogeneous case.
The surviving couplings are:
the expansion ($\Theta$), which drives the adiabatic
redshift of each species;
the shear ($\sigma_{ab}$), which couples
$\ell \leftrightarrow \ell \pm 2$ and is the primary
source of the CMB quadrupole; and
the vorticity ($\omega_{ab}$) and acceleration
($\dot{u}_a$), which provide fixed-$\ell$ mixing and
dipole sources respectively.

For an orthogonal Bianchi model
($\omega_{ab} = 0$, $\dot{u}_a = 0$), the
system~\eqref{eq:projected-system} simplifies to
\begin{equation}\label{eq:Bianchi-orth}
J_s(\alpha_s)\, \dot{\alpha}_s
= \mathcal{H}_s[\Theta;\, \alpha_s]
+ \mathcal{S}_s[\sigma_{ab};\, \alpha_s]
+ \mathcal{C}_s\,,
\end{equation}
where $\mathcal{H}_s$ collects the expansion-driven
redshift terms and $\mathcal{S}_s$ collects the
shear-driven coupling terms.  The structure of
$\mathcal{S}_s$ has a decisive consequence: starting
from an isotropic monopole state
($\Theta^{(s)}_{A_\ell} = 0$ for $\ell \geq 1$), the
shear $\sigma_{ab}$ sources the quadrupole
$\Theta^{(s)}_{ab}$ directly, while the dipole
$\Theta^{(s)}_a$ receives no shear source in the
orthogonal case.  The dipole is sourced only by tilt
(species drift relative to $n^a$), observer motion, or
inter-species momentum exchange.  This asymmetry
explains why the CMB quadrupole is the natural
observable for constraining the background shear, while
the dipole constrains the tilt sector---a separation
that is central to the departure decomposition
(Section~\ref{sec:master-id}).


\paragraph{Minimal truncation.}

The $(L_\Theta, L_\eta) = (2, 1)$ truncation retains the
temperature field through the quadrupole and the fugacity
field through the dipole:
\begin{equation}\label{eq:minimal-state}
\alpha_s
= \big(
\ln T_{0,s},\;
\eta_{0,s},\;
\Theta^{(s)}_{\langle a \rangle},\;
\Theta^{(s)}_{\langle ab \rangle},\;
\eta^{(s)}_{\langle a \rangle}
\big)\,.
\end{equation}
This truncation is the covariant lift of the $(E_0, E_1, E_2, N_0,
N_1)$ minimal observable set identified in the
slab-symmetric internal analysis: the energy sector alone
($E_0, E_1, E_2$)
cannot determine $\eta_{0,s}$ or $\eta^{(s)}_a$, and
omitting the number moments $N_0$ and $N_1$ renders the
fugacity sector unidentifiable.  The total number of
independent components for a single two-field species is
$1 + 1 + 3 + 5 + 3 = 13$; for a one-field species, it
reduces to $1 + 3 + 5 = 9$.  Summing over
$\mathcal{K} = \{\gamma, \nu_e, \bar{\nu}_e, \nu_x,
\bar{\nu}_x\}$ with one-field photons and two-field
neutrinos gives a total state-vector dimension of
$9 + 4 \times 13 = 61$ per Bianchi evaluation---a
number that is large but tractable for modern nested
sampling.


% ============================================================
\subsection{Limit recovery and scope delineation}
\label{sec:limit-recovery}
% ============================================================

We now verify that the multi-fluid framework
reduces correctly to all known special cases and
state explicitly the boundaries of the present
formalism.


\paragraph{FLRW isotropic limit.}

Setting $\Theta_s(\hat{e}) = 1$,
$\eta_s(\hat{e}) = \eta_{0,s}$, and all Bianchi
kinematic quantities to zero, the
system~\eqref{eq:projected-system} reduces to the
$2 \times 2$ monopole
system~\eqref{eq:FLRW-monopole}, and the departure
parameter $x_C = 0$ by the master identity.  This
limit is verified in Section~\ref{sec:species-conservation}
and matches the standard multi-component Friedmann
cosmology.


\paragraph{Single-species limit.}

Setting $|\mathcal{K}| = 1$ (a single kinetic species
with no fluid sector), the interaction sources vanish
($Q_s^b = 0$, $R_s = 0$), the total stress-energy is
$T^{ab}_{\rm tot} = T_s^{ab}$, and the projected
system~\eqref{eq:projected-system} reduces to the
single-species mass-matrix ODE of
Section~\ref{sec:ode-system}.  The tangency diagnostic
$D_{s,\geq 2}$ reduces to the single-species
diagnostic of Section~\ref{sec:tangency}.  The
$\Theta^4$ moment map
(Section~\ref{sec:moment-map}) and the nonlinear
correction factors $R_\sigma$, $R_\omega$, $R_{\dot{u}}$
(Section~\ref{sec:universal}) are recovered as the
$\eta_s = 0$, single-$s$ specialisation of the
multi-species reconstruction.


\paragraph{Scalar-$\eta$ limit.}

For a two-field species with
$\eta_s(\hat{e}) = \eta_{0,s}$ (uniform fugacity,
no angular structure in $\eta$), the number drift
$j_s^a = 0$ (Eq.~\ref{eq:j-s-recon}), and the
number sector decouples from the angular dynamics.
The species behaves as a one-field species with a
modified thermodynamic background---the integrals
$\mathcal{I}_n(\xi_s, \eta_{0,s})$ replace the
$\eta = 0$ values, but the angular coupling
structure is unchanged.  This limit is relevant for
neutrinos with a frozen homogeneous chemical potential
after decoupling.


\paragraph{Bianchi $\to$ FLRW limit of the hierarchy.}

Setting $\sigma_{ab} \to 0$ in the Bianchi
system~\eqref{eq:Bianchi-orth} removes the
$\ell \leftrightarrow \ell \pm 2$ shear coupling.
The quadrupole $\Theta^{(s)}_{ab}$ then evolves
under collisional damping alone and decays to zero,
while the monopole evolution reverts to the FLRW
system~\eqref{eq:FLRW-monopole}.  The departure
parameter $x_C \to 0$ continuously, confirming that
the multi-species framework does not introduce
spurious anisotropy in the isotropic limit---a
non-trivial check given the additional degrees of
freedom.


\paragraph{Scope: what is established and what is not.}

We close this section with a precise delineation of
the status of each component of the multi-fluid
framework, using the three-tier claim taxonomy
adopted throughout this work.

\medskip\noindent\textsc{Established.}\enspace
The common-frame contract
(Remark~\ref{rem:frame-contract}) and species
classification (Definition~\ref{def:sectors}) are
structural choices that do not depend on
approximations.
The macroscopic reconstruction formulae
(Eqs.~\ref{eq:rho-s-recon}--\ref{eq:j-s-recon})
are exact integrals of the $\Teff$ ansatz, and the
conservation and Einstein coupling
(Eqs.~\ref{eq:total-conservation}--\ref{eq:Einstein-multi})
are exact consequences of general relativity.
The Gram-matrix identifiability
(Eq.~\ref{eq:Gram-matrix}) is a structural property
of the $\Teff$ manifold geometry, and the tangency
diagnostic $D_{s,\geq 2}$
(Eq.~\ref{eq:D-diagnostic}) is a well-defined
per-species, per-ray functional of the collision
integral.
All four limit recoveries---FLRW, single-species,
scalar-$\eta$, and Bianchi$\to$FLRW---are verified
analytically.

\medskip\noindent\textsc{Conditional.}\enspace
The projected evolution
system~\eqref{eq:projected-system} is structurally
exact, but its explicit coefficient tables
(the entries of $J_s$, $\mathcal{R}_s$,
$\mathcal{C}_s$) are derived here only in the
isotropic limit (Eq.~\ref{eq:FLRW-monopole}) and
the linearised Bianchi limit.  The full nonlinear
Jacobian for arbitrary $(L_\Theta, L_\eta)$ requires
Gaunt-type angular integrals that are deferred to the
numerical implementation
(Chapter~\ref{ch:results}).
The minimal truncation
$(L_\Theta, L_\eta) = (2, 1)$ is sufficient for
the scalar-amplitude pipeline at the current
precision; whether it remains sufficient for the
direction-dependent programme
(Section~\ref{sec:open-problems}) depends on the
magnitude of $D_{s,\geq 2}$ at the required
sky resolution.

\medskip\noindent\textsc{Not established.}\enspace
Four items remain open.
First, global-in-time existence and stability of the
coupled multi-species system: the rigorous
well-posedness result of Fournodavlos, Marshall \&
Oliynyk~\cite{FournodavlosMarshallOliynyk2025}
covers only the two-fluid super-radiative case
($1/3 < c_s^2 < 3/7$), and the physically relevant
dust--radiation--$\Lambda$ system remains an open
problem (Section~\ref{sec:open-problems}).
Second, symmetric hyperbolicity of the projected
system~\eqref{eq:projected-system}.
Third, robustness of the $\Teff$ reduction under
realistic microphysical collision kernels beyond
Thomson scattering.
Fourth, the equivalence between the local adequacy
certificate $D_{s,\geq 2} < \varepsilon_{\rm thr}$
and dynamical adequacy over cosmological
timescales---the tangency diagnostic is a
necessary but not sufficient condition for
long-time accuracy.

These open items define the boundary between the
results of this work and the programme of future
extensions discussed in
Chapter~\ref{ch:future}.

% ============================================================
\section{Perturbation sector: harmonic modes on Bianchi backgrounds}
\label{sec:perturbation-bianchi-modes}
% ============================================================

The background-only Lowell system closes the homogeneous Bianchi
sector, but the observable CMB anisotropy requires the
inhomogeneous perturbation problem at $k \neq 0$.  In the present
implementation we factor that extension into seven explicit
surfaces: a harmonic-mode context builder, a rotated
$\tilde{\nabla}$ operator, a regular adiabatic initial-condition
constructor, a $k \rightarrow 0$ recovery gate, a Class-B
quantisation helper, a tilted-seed boost rule, and a
type-by-type regression harness.  This decomposition keeps the
homogeneous LB-6 background anchor intact while making every
perturbative assumption auditable in code and in figures.

\subsection{Harmonic-mode decomposition and complex transport}

The perturbation state is labelled by a harmonic mode
$Y_k(x)$ adapted to the relevant Bianchi type.  Types~I and
FLRW reduce to ordinary plane waves, Type~V carries the
hyperbolic offset of Harrison's open-universe harmonics
\cite{Harrison1967}, Type~VII$_0$ keeps the symmetric-line helical
reduction, and Type~IX retains the discrete $S^3$ spectrum.
Class-B types (III, IV, VI$_h$, VII$_h$) are represented on the
abelian $(e_1,e_3)$ plane, with the spiral morphology of
VII$_h$ following the Pontzen--Challinor description
\cite{PontzenChallinor2007}.  The crucial bookkeeping point is
that $\tilde{\nabla}_a$ acting on a mode is generically complex,
so the hierarchy driver must preserve complex dtype all the way
through the packed PSTF tower.

\subsection{Full $\tilde{\nabla}$ operator}

Axis-aligned subsets are already covered by the audited
FB-2 dispatch.  The FB-5 lift adds a full rotated operator by
constructing the verified axis-aligned operator in a canonical
frame and conjugating it with an explicit Euler rotation.  This
makes the off-axis Wigner-style dependence visible at the API
level rather than burying it in hidden control flow, and it
ensures that the identity-angle limit reduces exactly to the
audited FB-2 branch.

\subsection{Regular adiabatic seed}

The regular adiabatic startup is anchored to the radiation-era
series for $\eta_{\rm cov}$, $\Delta_\gamma$, $\Delta_b$,
$\Delta_c$, $\Delta_\nu$, $q_\gamma$, $q_\nu$, $\pi_\nu$, and
$G_3$, with the photon quadrupole placed on the tight-coupling
manifold rather than forced to zero.  In the flat FLRW limit this
is the same superhorizon regular mode that underlies the standard
CAMB/CLASS startup and the Sachs--Wolfe large-scale picture
\cite{SachsWolfe1967}.  the conditioned legacy figure gallery in Appendix~\ref{app:conditioned-legacy-figure-gallery}
summarises the per-type harmonic families, while
the conditioned legacy figure gallery in Appendix~\ref{app:conditioned-legacy-figure-gallery} shows the startup
scalings for the photon and matter perturbations.

\subsection{$k = 0$ limit gate}

The $k = 0$ branch is not a separate solver; it is a recovery
condition on the perturbative pipeline.  We therefore expose it as
an explicit validator that compares the perturbative prefix
against the LB-6 zero-state anchor and requires byte identity at
exact $k = 0$.  This sharpens two separate claims: the harmonic
layer may introduce complex transport at $k \neq 0$, but it must
collapse continuously onto the background-only branch as the mode
scale vanishes; and any extra perturbation auxiliaries must also
vanish in that limit.

\subsection{Class-B quantisation}

The Class-B families require a mode label beyond the raw
three-vector because the scalar Laplacian carries a twist offset.
For Type~V the offset is Harrison's $a^2$ term
\cite{Harrison1967}; for III, IV, VI$_h$, and VII$_h$ the
implementation generalises this to $a_{\rm twist}^2/(1+|h|)$
using the same structure-constant single source of truth as the
background sector.  The result is a metadata layer that can be
threaded through the perturbation pipeline without widening the
core harmonic descriptor.

\subsection{Tilted seed regularisation}

Tilted perturbation startup is treated as a two-step map: first
construct the regular seed in the electron frame, then apply the
axisymmetric boost kernel and re-project onto the retained
$m=0$ surface.  The explicit separation matters because the
$\beta = 0$ path must remain byte-identical to the orthogonal
seed, while the non-zero tilt case mixes adjacent multipoles in
the way anticipated by the PSTF boost algebra.

\subsection{$k \times$ type regression surface}

The final perturbation layer is a regression harness over
Bianchi type and wavenumber.  Its role is not to hide the oracle
inside an opaque test, but to make the coverage grid explicit:
type label, mode family, spectrum kind, regular seed, and
low-$\ell$ $D_\ell^{TT}$ comparison are all returned in one
auditable bundle.  In the FLRW and Type-I limit the harness is
compared directly to the Planck-2018 CAMB reference, while the
other types retain their geometry-dependent modulation on the
same shared $\ell$ grid.

% === END INLINED ===

% === INLINED: figures_ch02.tex ===
% Figures for Chapter 2
% Updated: TF-D06, 2026-03-12
% Added: fig_tilted_dictionary, fig_peculiar_jeans, fig_vorticity_hierarchy

% fig_peculiar_jeans placed in ch09



% === END INLINED ===

% BEGIN current-code generated figure deck
% Generated by scripts/make_current_manuscript_figures.py.
% Do not edit figure paths by hand; regenerate instead.

\begin{figure}[htbp]
\centering
\includegraphics[width=0.92\textwidth]{current/fig_current_dag_progress}
\caption{Current DAG/status snapshot by repository owner. The panel is bookkeeping evidence only: completion and smoke status do not promote any result to production validation.}
\label{fig:current-dag-progress}
\end{figure}

\begin{figure}[htbp]
\centering
\includegraphics[width=0.92\textwidth]{current/fig_current_scalar_morphology_boundary}
\caption{Current scalar-to-morphology comparison matrix. The allowed state is diagnostic side-by-side reporting with separate provenance; scalar diagnostics and low-ell summaries do not identify geometry or provide Bianchi family-ID.}
\label{fig:current-scalar-morphology-boundary}
\end{figure}

% END current-code generated figure deck
